%% ============================================================================
%%  Elastic Rod Manipulation Using a Contact Point
%%  RSS 2027 submission -- main file (preamble, macros, title, abstract, sections)
%%
%%  Template: official RSS LaTeX template (IEEEtran.cls v1.7a + natbib/plainnat),
%%  roboticsconference.org/docs/paper-template-latex.tar.gz, fetched 2026-11-24.
%%  Do NOT change fonts, margins, spacing, or IEEEtran.cls: the CFP may
%%  desk-reject submissions that alter the template formatting.
%%
%%  RSS 2027 two-stage process (roboticsconference.org/information/cfp/):
%%   Stage 1  Extended Abstract, due Dec 4, 2026, 11:59pm AoE (OpenReview).
%%            Max 5 pages of content (appendices count) + 1 page of references.
%%            Problem statement, testable hypotheses, and methodology must be
%%            fully developed; experiments may be missing or preliminary.
%%            Related work: only the most direct citations. No supplementary
%%            material; websites and appendices are not reviewed.
%%   Stage 2  Final paper, due Apr 16, 2027 (if invited): max 8 pages excluding
%%            references; executes the Extension Charter; complete related work
%%            and experiments. Supplementary material (zip, 100MB) Apr 19, 2027.
%%  Double-blind: no names, affiliations, acknowledgments, or links (not even
%%  anonymous ones); anonymized PDF metadata; self-citations in third person;
%%  no lab logos in photos.
%%  Also: no page numbers, no "Draft"/"To appear" marks, Type-1 fonts only,
%%  and no hallucinated references (checked automatically -> desk reject).
%%
%%  Build (draft, -> main.pdf):
%%    latexmk -pdf main.tex
%%  Build (anonymous, notes hidden, -> main_submission.pdf):
%%    latexmk -pdf -jobname=main_submission -usepretex='\def\papermode{submission}' main.tex
%%  With the Makefile: make | make submission | make check (pages, fonts, metadata)
%% ============================================================================
\documentclass[conference]{IEEEtran}

%% ---------------------------------------------------------------------------
%% Build mode: draft | submission | camera
%%   draft       author names + TODO/author notes shown (default; this adds no
%%               "Draft" marking to the PDF)
%%   submission  anonymous author block and PDF metadata, notes hidden
%%               (use for both RSS review stages)
%%   camera      author names shown, notes hidden (camera-ready)
%% Change the default here or override from the command line, e.g.
%%   pdflatex "\def\papermode{submission}%% ============================================================================
%%  Elastic Rod Manipulation Using a Contact Point
%%  RSS 2027 submission -- main file (preamble, macros, title, abstract, sections)
%%
%%  Template: official RSS LaTeX template (IEEEtran.cls v1.7a + natbib/plainnat),
%%  roboticsconference.org/docs/paper-template-latex.tar.gz, fetched 2026-11-24.
%%  Do NOT change fonts, margins, spacing, or IEEEtran.cls: the CFP may
%%  desk-reject submissions that alter the template formatting.
%%
%%  RSS 2027 two-stage process (roboticsconference.org/information/cfp/):
%%   Stage 1  Extended Abstract, due Dec 4, 2026, 11:59pm AoE (OpenReview).
%%            Max 5 pages of content (appendices count) + 1 page of references.
%%            Problem statement, testable hypotheses, and methodology must be
%%            fully developed; experiments may be missing or preliminary.
%%            Related work: only the most direct citations. No supplementary
%%            material; websites and appendices are not reviewed.
%%   Stage 2  Final paper, due Apr 16, 2027 (if invited): max 8 pages excluding
%%            references; executes the Extension Charter; complete related work
%%            and experiments. Supplementary material (zip, 100MB) Apr 19, 2027.
%%  Double-blind: no names, affiliations, acknowledgments, or links (not even
%%  anonymous ones); anonymized PDF metadata; self-citations in third person;
%%  no lab logos in photos.
%%  Also: no page numbers, no "Draft"/"To appear" marks, Type-1 fonts only,
%%  and no hallucinated references (checked automatically -> desk reject).
%%
%%  Build (draft, -> main.pdf):
%%    latexmk -pdf main.tex
%%  Build (anonymous, notes hidden, -> main_submission.pdf):
%%    latexmk -pdf -jobname=main_submission -usepretex='\def\papermode{submission}' main.tex
%%  With the Makefile: make | make submission | make check (pages, fonts, metadata)
%% ============================================================================
\documentclass[conference]{IEEEtran}

%% ---------------------------------------------------------------------------
%% Build mode: draft | submission | camera
%%   draft       author names + TODO/author notes shown (default; this adds no
%%               "Draft" marking to the PDF)
%%   submission  anonymous author block and PDF metadata, notes hidden
%%               (use for both RSS review stages)
%%   camera      author names shown, notes hidden (camera-ready)
%% Change the default here or override from the command line, e.g.
%%   pdflatex "\def\papermode{submission}%% ============================================================================
%%  Elastic Rod Manipulation Using a Contact Point
%%  RSS 2027 submission -- main file (preamble, macros, title, abstract, sections)
%%
%%  Template: official RSS LaTeX template (IEEEtran.cls v1.7a + natbib/plainnat),
%%  roboticsconference.org/docs/paper-template-latex.tar.gz, fetched 2026-11-24.
%%  Do NOT change fonts, margins, spacing, or IEEEtran.cls: the CFP may
%%  desk-reject submissions that alter the template formatting.
%%
%%  RSS 2027 two-stage process (roboticsconference.org/information/cfp/):
%%   Stage 1  Extended Abstract, due Dec 4, 2026, 11:59pm AoE (OpenReview).
%%            Max 5 pages of content (appendices count) + 1 page of references.
%%            Problem statement, testable hypotheses, and methodology must be
%%            fully developed; experiments may be missing or preliminary.
%%            Related work: only the most direct citations. No supplementary
%%            material; websites and appendices are not reviewed.
%%   Stage 2  Final paper, due Apr 16, 2027 (if invited): max 8 pages excluding
%%            references; executes the Extension Charter; complete related work
%%            and experiments. Supplementary material (zip, 100MB) Apr 19, 2027.
%%  Double-blind: no names, affiliations, acknowledgments, or links (not even
%%  anonymous ones); anonymized PDF metadata; self-citations in third person;
%%  no lab logos in photos.
%%  Also: no page numbers, no "Draft"/"To appear" marks, Type-1 fonts only,
%%  and no hallucinated references (checked automatically -> desk reject).
%%
%%  Build (draft, -> main.pdf):
%%    latexmk -pdf main.tex
%%  Build (anonymous, notes hidden, -> main_submission.pdf):
%%    latexmk -pdf -jobname=main_submission -usepretex='\def\papermode{submission}' main.tex
%%  With the Makefile: make | make submission | make check (pages, fonts, metadata)
%% ============================================================================
\documentclass[conference]{IEEEtran}

%% ---------------------------------------------------------------------------
%% Build mode: draft | submission | camera
%%   draft       author names + TODO/author notes shown (default; this adds no
%%               "Draft" marking to the PDF)
%%   submission  anonymous author block and PDF metadata, notes hidden
%%               (use for both RSS review stages)
%%   camera      author names shown, notes hidden (camera-ready)
%% Change the default here or override from the command line, e.g.
%%   pdflatex "\def\papermode{submission}%% ============================================================================
%%  Elastic Rod Manipulation Using a Contact Point
%%  RSS 2027 submission -- main file (preamble, macros, title, abstract, sections)
%%
%%  Template: official RSS LaTeX template (IEEEtran.cls v1.7a + natbib/plainnat),
%%  roboticsconference.org/docs/paper-template-latex.tar.gz, fetched 2026-11-24.
%%  Do NOT change fonts, margins, spacing, or IEEEtran.cls: the CFP may
%%  desk-reject submissions that alter the template formatting.
%%
%%  RSS 2027 two-stage process (roboticsconference.org/information/cfp/):
%%   Stage 1  Extended Abstract, due Dec 4, 2026, 11:59pm AoE (OpenReview).
%%            Max 5 pages of content (appendices count) + 1 page of references.
%%            Problem statement, testable hypotheses, and methodology must be
%%            fully developed; experiments may be missing or preliminary.
%%            Related work: only the most direct citations. No supplementary
%%            material; websites and appendices are not reviewed.
%%   Stage 2  Final paper, due Apr 16, 2027 (if invited): max 8 pages excluding
%%            references; executes the Extension Charter; complete related work
%%            and experiments. Supplementary material (zip, 100MB) Apr 19, 2027.
%%  Double-blind: no names, affiliations, acknowledgments, or links (not even
%%  anonymous ones); anonymized PDF metadata; self-citations in third person;
%%  no lab logos in photos.
%%  Also: no page numbers, no "Draft"/"To appear" marks, Type-1 fonts only,
%%  and no hallucinated references (checked automatically -> desk reject).
%%
%%  Build (draft, -> main.pdf):
%%    latexmk -pdf main.tex
%%  Build (anonymous, notes hidden, -> main_submission.pdf):
%%    latexmk -pdf -jobname=main_submission -usepretex='\def\papermode{submission}' main.tex
%%  With the Makefile: make | make submission | make check (pages, fonts, metadata)
%% ============================================================================
\documentclass[conference]{IEEEtran}

%% ---------------------------------------------------------------------------
%% Build mode: draft | submission | camera
%%   draft       author names + TODO/author notes shown (default; this adds no
%%               "Draft" marking to the PDF)
%%   submission  anonymous author block and PDF metadata, notes hidden
%%               (use for both RSS review stages)
%%   camera      author names shown, notes hidden (camera-ready)
%% Change the default here or override from the command line, e.g.
%%   pdflatex "\def\papermode{submission}\input{main}"
%% ---------------------------------------------------------------------------
\providecommand*{\papermode}{draft}
\newcommand*{\paperid}{[add your ID here]} % OpenReview paper ID (Stage 1)
\newcommand*{\papertitle}{Elastic Rod Manipulation Using a Contact Point}

%% ---------------------------------------------------------------------------
%% Packages from the RSS template (hyperref is loaded further down, after the
%% packages it must follow)
%% ---------------------------------------------------------------------------
\usepackage{times}
% numbers option provides compact numerical references in the text.
\usepackage[numbers]{natbib}
\usepackage{multicol}

%% ---------------------------------------------------------------------------
%% Additional packages (none of them changes fonts, margins, or spacing)
%% Avoid: subcaption/caption, enumitem, geometry, mathptmx/newtx, titlesec --
%% they alter the IEEEtran layout.
%% ---------------------------------------------------------------------------
\usepackage{amsmath}
\usepackage{amssymb}
\usepackage{bm}
\usepackage{graphicx}
\graphicspath{{figures/}}
\usepackage[caption=false,font=footnotesize]{subfig} % IEEEtran-safe subfigures
\usepackage{booktabs}
\usepackage{multirow}
\usepackage{siunitx}  % \SI{2}{\milli\metre}, \SI{5}{\degree}
\usepackage{xcolor}
\usepackage{algorithm}
\usepackage{algpseudocode}
\usepackage{xspace}
\usepackage{etoolbox}

\usepackage[bookmarks=true]{hyperref} % RSS template package
\usepackage[capitalise]{cleveref}     % must come after hyperref
\crefname{section}{Sec.}{Secs.}
\crefname{subsection}{Sec.}{Secs.}
\crefname{figure}{Fig.}{Figs.}
\crefname{table}{Table}{Tables}
\crefname{equation}{Eq.}{Eqs.}
\crefname{algorithm}{Alg.}{Algs.}

%% ---------------------------------------------------------------------------
%% Build-mode switches
%% ---------------------------------------------------------------------------
\newif\ifanonymous
\newif\ifshownotes
\ifdefstring{\papermode}{submission}{\anonymoustrue\shownotesfalse}{%
  \ifdefstring{\papermode}{camera}{\anonymousfalse\shownotesfalse}{%
    \anonymousfalse\shownotestrue}}% draft (also the fallback)
\typeout{*** Paper mode: \papermode}

%% PDF metadata (must be anonymous in the review stages)
\hypersetup{
  pdftitle    = {\papertitle},
  pdfkeywords = {deformable linear objects; elastic rods; contact; haptic sensing; manipulation},
}
\ifanonymous
  \hypersetup{pdfauthor={}}
\else
  \hypersetup{pdfauthor={Stav Ashur, Yuval Cassuto, Avishai Sintov}}
\fi

%% ---------------------------------------------------------------------------
%% Draft notes: visible only in draft mode, removed in submission/camera modes
%% (do not also load the todonotes package: it defines \todo as well)
%% ---------------------------------------------------------------------------
\definecolor{notecolor}{RGB}{200,0,0}
\definecolor{outlinecolor}{RGB}{0,70,150}
\DeclareRobustCommand{\draftnote}[3]{%
  \ifshownotes\textcolor{#1}{[\textbf{#2:} #3]}\else\ifhmode\unskip\fi\fi}
\newcommand{\todo}[1]{\draftnote{notecolor}{TODO}{#1}}
\newcommand{\stav}[1]{\draftnote{orange}{Stav}{#1}}
\newcommand{\yuval}[1]{\draftnote{teal}{Yuval}{#1}}
\newcommand{\avishai}[1]{\draftnote{violet}{Avishai}{#1}}

%% Skeleton outlines: bullet lists of planned content. Replace each outline
%% block with prose as you write; they stay blue in every mode so leftovers
%% are easy to spot, and non-draft builds log how many are left.
\newcounter{outlineblocks}
\newenvironment{outline}{%
  \stepcounter{outlineblocks}\par\color{outlinecolor}\begin{itemize}}%
  {\end{itemize}}
\AtEndDocument{\ifshownotes\else\ifnum\value{outlineblocks}>0
  \typeout{WARNING: \arabic{outlineblocks} skeleton outline block(s) left in the paper}%
\fi\fi}

%% \figplaceholder{<width>}{<height>}{<description>}: framed stand-in for a
%% figure that does not exist yet. Replace with \includegraphics{...}.
\newcommand{\figplaceholder}[3]{%
  \fbox{\parbox[c][#2][c]{\dimexpr #1-2\fboxsep-2\fboxrule\relax}{%
    \centering\footnotesize\color{outlinecolor}#3}}}

%% ---------------------------------------------------------------------------
%% Hypotheses and problem statements
%%   \begin{hypothesis}[Short name]\label{hyp:x} ... \end{hypothesis}
%%   \ref{hyp:x} prints "H1", "H2", ...
%% ---------------------------------------------------------------------------
\newcounter{hypothesis}
\renewcommand{\thehypothesis}{H\arabic{hypothesis}}
\crefname{hypothesis}{hypothesis}{hypotheses}
\newenvironment{hypothesis}[1][]{%
  \par\smallskip\refstepcounter{hypothesis}\noindent
  \textbf{\thehypothesis\ifstrempty{#1}{}{ (#1)}:}\hspace{0.5em}\itshape\ignorespaces}%
  {\par\smallskip}
\newtheorem{problem}{Problem}
\crefname{problem}{Problem}{Problems}

%% ---------------------------------------------------------------------------
%% Notation
%% ---------------------------------------------------------------------------
\newcommand{\R}{\mathbb{R}}
\newcommand{\SE}[1]{\mathrm{SE}(#1)}
\newcommand{\SO}[1]{\mathrm{SO}(#1)}
\newcommand{\vect}[1]{\bm{#1}}
\newcommand{\gpose}{\vect{T}_{\mathrm{g}}}          % gripper (end-effector) pose
\newcommand{\tpose}{\vect{T}_{\mathrm{tip}}}        % DLO tip pose
\newcommand{\tposehat}{\hat{\vect{T}}_{\mathrm{tip}}} % estimated tip pose
\newcommand{\tposegoal}{\vect{T}^{*}_{\mathrm{tip}}} % target tip pose
\newcommand{\wrench}{\vect{w}}                      % wrist F/T reading (f, tau)
\newcommand{\cpt}{\vect{c}}                         % contact measurement
\newcommand{\cang}{\phi_{\mathrm{c}}}               % contact angle on the ring
\newcommand{\Rfree}{\mathcal{R}_{\mathrm{free}}}    % reachable set, no contact
\newcommand{\Rcon}{\mathcal{R}_{\mathrm{c}}}        % reachable set, with contact

\newcommand{\eg}{e.g.,\xspace}
\newcommand{\ie}{i.e.,\xspace}
\newcommand{\etal}{et~al.\xspace}
\newcommand{\wrt}{w.r.t.\xspace}
\newcommand{\cf}{cf.\xspace}

\begin{document}

\title{\papertitle}

\ifanonymous
  % Author block required by the RSS template for anonymous review
  \author{Author Names Omitted for Anonymous Review. Paper-ID \paperid}
\else
  \author{\authorblockN{Stav Ashur, Yuval Cassuto, and Avishai Sintov}
  \authorblockA{Tel Aviv University, Tel Aviv, Israel}}
  % TODO (camera-ready): departments and e-mail addresses
\fi

\maketitle

\begin{abstract}
Sensors and devices such as cameras and brushes can be mounted on the tips of
elastic rods and other deformable linear objects (DLOs), allowing them to be
inserted into and manipulated within constrained spaces. Using the right
contact point enables the DLO to reach poses that are infeasible when no
contact with the environment is acceptable. Such a contact point serves as an
ad-hoc joint, facilitating direction changes of the DLO's tip that are
impossible if no collisions are tolerated. In this work, we propose the use of
haptic sensors to perform manipulation that brings the tip of the DLO to a
target pose in a constrained environment by utilizing a contact point of the
DLO with an environment edge.
\todo{Add 1--2 sentences on the hypotheses and methodology (vision-free tip-pose
estimation from gripper pose, wrist force/torque, and contact location;
contact-aware steering) and on the preliminary results.}
\end{abstract}

\IEEEpeerreviewmaketitle

%% Sections, in order
\input{sections/intro}
\input{sections/related_work}
\input{sections/problem_statement}
\input{sections/hypotheses}
\input{sections/method}
\input{sections/experiments}
\input{sections/conclusion}

%% Acknowledgments are not allowed in the anonymous submission.
\ifanonymous\else
\section*{Acknowledgments}
\todo{Funding and acknowledgments (camera-ready only).}
\fi

%% Use plainnat to work nicely with natbib (RSS template).
\bibliographystyle{plainnat}
\bibliography{references}

\end{document}
"
%% ---------------------------------------------------------------------------
\providecommand*{\papermode}{draft}
\newcommand*{\paperid}{[add your ID here]} % OpenReview paper ID (Stage 1)
\newcommand*{\papertitle}{Elastic Rod Manipulation Using a Contact Point}

%% ---------------------------------------------------------------------------
%% Packages from the RSS template (hyperref is loaded further down, after the
%% packages it must follow)
%% ---------------------------------------------------------------------------
\usepackage{times}
% numbers option provides compact numerical references in the text.
\usepackage[numbers]{natbib}
\usepackage{multicol}

%% ---------------------------------------------------------------------------
%% Additional packages (none of them changes fonts, margins, or spacing)
%% Avoid: subcaption/caption, enumitem, geometry, mathptmx/newtx, titlesec --
%% they alter the IEEEtran layout.
%% ---------------------------------------------------------------------------
\usepackage{amsmath}
\usepackage{amssymb}
\usepackage{bm}
\usepackage{graphicx}
\graphicspath{{figures/}}
\usepackage[caption=false,font=footnotesize]{subfig} % IEEEtran-safe subfigures
\usepackage{booktabs}
\usepackage{multirow}
\usepackage{siunitx}  % \SI{2}{\milli\metre}, \SI{5}{\degree}
\usepackage{xcolor}
\usepackage{algorithm}
\usepackage{algpseudocode}
\usepackage{xspace}
\usepackage{etoolbox}

\usepackage[bookmarks=true]{hyperref} % RSS template package
\usepackage[capitalise]{cleveref}     % must come after hyperref
\crefname{section}{Sec.}{Secs.}
\crefname{subsection}{Sec.}{Secs.}
\crefname{figure}{Fig.}{Figs.}
\crefname{table}{Table}{Tables}
\crefname{equation}{Eq.}{Eqs.}
\crefname{algorithm}{Alg.}{Algs.}

%% ---------------------------------------------------------------------------
%% Build-mode switches
%% ---------------------------------------------------------------------------
\newif\ifanonymous
\newif\ifshownotes
\ifdefstring{\papermode}{submission}{\anonymoustrue\shownotesfalse}{%
  \ifdefstring{\papermode}{camera}{\anonymousfalse\shownotesfalse}{%
    \anonymousfalse\shownotestrue}}% draft (also the fallback)
\typeout{*** Paper mode: \papermode}

%% PDF metadata (must be anonymous in the review stages)
\hypersetup{
  pdftitle    = {\papertitle},
  pdfkeywords = {deformable linear objects; elastic rods; contact; haptic sensing; manipulation},
}
\ifanonymous
  \hypersetup{pdfauthor={}}
\else
  \hypersetup{pdfauthor={Stav Ashur, Yuval Cassuto, Avishai Sintov}}
\fi

%% ---------------------------------------------------------------------------
%% Draft notes: visible only in draft mode, removed in submission/camera modes
%% (do not also load the todonotes package: it defines \todo as well)
%% ---------------------------------------------------------------------------
\definecolor{notecolor}{RGB}{200,0,0}
\definecolor{outlinecolor}{RGB}{0,70,150}
\DeclareRobustCommand{\draftnote}[3]{%
  \ifshownotes\textcolor{#1}{[\textbf{#2:} #3]}\else\ifhmode\unskip\fi\fi}
\newcommand{\todo}[1]{\draftnote{notecolor}{TODO}{#1}}
\newcommand{\stav}[1]{\draftnote{orange}{Stav}{#1}}
\newcommand{\yuval}[1]{\draftnote{teal}{Yuval}{#1}}
\newcommand{\avishai}[1]{\draftnote{violet}{Avishai}{#1}}

%% Skeleton outlines: bullet lists of planned content. Replace each outline
%% block with prose as you write; they stay blue in every mode so leftovers
%% are easy to spot, and non-draft builds log how many are left.
\newcounter{outlineblocks}
\newenvironment{outline}{%
  \stepcounter{outlineblocks}\par\color{outlinecolor}\begin{itemize}}%
  {\end{itemize}}
\AtEndDocument{\ifshownotes\else\ifnum\value{outlineblocks}>0
  \typeout{WARNING: \arabic{outlineblocks} skeleton outline block(s) left in the paper}%
\fi\fi}

%% \figplaceholder{<width>}{<height>}{<description>}: framed stand-in for a
%% figure that does not exist yet. Replace with \includegraphics{...}.
\newcommand{\figplaceholder}[3]{%
  \fbox{\parbox[c][#2][c]{\dimexpr #1-2\fboxsep-2\fboxrule\relax}{%
    \centering\footnotesize\color{outlinecolor}#3}}}

%% ---------------------------------------------------------------------------
%% Hypotheses and problem statements
%%   \begin{hypothesis}[Short name]\label{hyp:x} ... \end{hypothesis}
%%   \ref{hyp:x} prints "H1", "H2", ...
%% ---------------------------------------------------------------------------
\newcounter{hypothesis}
\renewcommand{\thehypothesis}{H\arabic{hypothesis}}
\crefname{hypothesis}{hypothesis}{hypotheses}
\newenvironment{hypothesis}[1][]{%
  \par\smallskip\refstepcounter{hypothesis}\noindent
  \textbf{\thehypothesis\ifstrempty{#1}{}{ (#1)}:}\hspace{0.5em}\itshape\ignorespaces}%
  {\par\smallskip}
\newtheorem{problem}{Problem}
\crefname{problem}{Problem}{Problems}

%% ---------------------------------------------------------------------------
%% Notation
%% ---------------------------------------------------------------------------
\newcommand{\R}{\mathbb{R}}
\newcommand{\SE}[1]{\mathrm{SE}(#1)}
\newcommand{\SO}[1]{\mathrm{SO}(#1)}
\newcommand{\vect}[1]{\bm{#1}}
\newcommand{\gpose}{\vect{T}_{\mathrm{g}}}          % gripper (end-effector) pose
\newcommand{\tpose}{\vect{T}_{\mathrm{tip}}}        % DLO tip pose
\newcommand{\tposehat}{\hat{\vect{T}}_{\mathrm{tip}}} % estimated tip pose
\newcommand{\tposegoal}{\vect{T}^{*}_{\mathrm{tip}}} % target tip pose
\newcommand{\wrench}{\vect{w}}                      % wrist F/T reading (f, tau)
\newcommand{\cpt}{\vect{c}}                         % contact measurement
\newcommand{\cang}{\phi_{\mathrm{c}}}               % contact angle on the ring
\newcommand{\Rfree}{\mathcal{R}_{\mathrm{free}}}    % reachable set, no contact
\newcommand{\Rcon}{\mathcal{R}_{\mathrm{c}}}        % reachable set, with contact

\newcommand{\eg}{e.g.,\xspace}
\newcommand{\ie}{i.e.,\xspace}
\newcommand{\etal}{et~al.\xspace}
\newcommand{\wrt}{w.r.t.\xspace}
\newcommand{\cf}{cf.\xspace}

\begin{document}

\title{\papertitle}

\ifanonymous
  % Author block required by the RSS template for anonymous review
  \author{Author Names Omitted for Anonymous Review. Paper-ID \paperid}
\else
  \author{\authorblockN{Stav Ashur, Yuval Cassuto, and Avishai Sintov}
  \authorblockA{Tel Aviv University, Tel Aviv, Israel}}
  % TODO (camera-ready): departments and e-mail addresses
\fi

\maketitle

\begin{abstract}
Sensors and devices such as cameras and brushes can be mounted on the tips of
elastic rods and other deformable linear objects (DLOs), allowing them to be
inserted into and manipulated within constrained spaces. Using the right
contact point enables the DLO to reach poses that are infeasible when no
contact with the environment is acceptable. Such a contact point serves as an
ad-hoc joint, facilitating direction changes of the DLO's tip that are
impossible if no collisions are tolerated. In this work, we propose the use of
haptic sensors to perform manipulation that brings the tip of the DLO to a
target pose in a constrained environment by utilizing a contact point of the
DLO with an environment edge.
\todo{Add 1--2 sentences on the hypotheses and methodology (vision-free tip-pose
estimation from gripper pose, wrist force/torque, and contact location;
contact-aware steering) and on the preliminary results.}
\end{abstract}

\IEEEpeerreviewmaketitle

%% Sections, in order
%% ============================================================================
%% I. Introduction                               target: ~0.75 page incl. Fig. 1
%% Motivate the problem, state the gap, preview the hypotheses (Sec. IV) and
%% the methodology (Sec. V), and list the contributions.
%% ============================================================================
\section{Introduction}
\label{sec:intro}

\begin{outline}
  \item \textbf{Motivation.} Cameras, brushes, and other tools are mounted on
    the tips of elastic rods and other DLOs and inserted into constrained
    spaces, \eg for inspection and repair inside engines~\cite{dong2017slender}.
    \todo{1--2 more application examples.}
  \item \textbf{Challenge.} The rod's shape depends on mechanics that are hard
    to model~\cite{sanchez2018survey}, and once inserted, the rod is hidden by
    the environment or by the robot itself, so visual state estimation is
    unreliable.
  \item \textbf{Key idea.} Contact with the environment is usually avoided.
    Instead, we use a contact point between the rod and an environment edge as
    an \emph{ad-hoc joint}: pivoting about it changes the tip direction in ways
    that are impossible if no collisions are tolerated (\cref{fig:teaser}).
  \item \textbf{Problem.} Bring the tip of a gripped elastic rod to a target
    pose in a constrained environment by exploiting a contact with an
    environment edge, using only haptic sensing (robot proprioception, wrist
    force/torque, contact location) during execution (\cref{sec:problem}).
  \item \textbf{Approach.} A sensorized edge reports where the rod touches it;
    a learned estimator maps gripper pose, wrench, and contact location to the
    tip pose; a contact-aware controller steers the tip to the target
    (\cref{sec:method}).
\end{outline}

% Fig. 1 is defined here, not at the top of the section, so that it lands at
% the top of the right column of page 1. LaTeX puts a [t] float at the top of
% the left column (above the abstract) if that column still has room when the
% float is reached; if this happens after edits, move this block further down.
\begin{figure}[t]
  \centering
  \figplaceholder{\linewidth}{4.2cm}{%
    \textbf{Fig.~1 (teaser).} A robot grasps an elastic rod whose tip carries a
    tool. The rod rests on an environment edge (sensorized ring) that acts as an
    ad-hoc joint, letting the tip turn toward a target pose that is unreachable
    without contact. Inset: reachable tip directions with vs.\ without contact.}
  \caption{Using a contact point with an environment edge as an ad-hoc joint
    lets the tip of a gripped elastic rod reach poses that are infeasible
    without contact.}
  \label{fig:teaser}
\end{figure}

\begin{outline}
  \item \textbf{Hypotheses} (\cref{sec:hypotheses}): contact enlarges the set
    of reachable tip poses (\ref{hyp:reach}); haptic signals suffice to estimate
    the tip pose without vision (\ref{hyp:estimate}); closed-loop haptic
    steering reaches targets that contact-free or open-loop strategies miss
    (\ref{hyp:steer}).
  \item \textbf{Contributions}:
    \begin{itemize}
      \item a formulation of DLO tip manipulation that uses a single
        environment contact as an ad-hoc joint, with a reachability analysis;
      \item a vision-free tip-pose estimator that uses proprioception,
        force/torque, and contact location;
      \item an experimental methodology and dataset with a Nitinol rod, from
        planar (2D) to spatial (3D) settings;
      \item preliminary results on the hypotheses.
        \todo{Update after the experiments.}
    \end{itemize}
\end{outline}

%% ============================================================================
%% II. Related Work                                           target: ~0.5 page
%% Stage 1: cite only the work most directly tied to the central problem (what
%% is needed to establish novelty). Stage 2 requires a complete section.
%% Anonymity: cite the authors' own prior work in the third person.
%% RSS 2027 checks references automatically and desk-rejects hallucinated
%% citations, so check every entry you add to references.bib.
%% Merged with the notes in material/2_Related_work.tex (earlier document).
%% ============================================================================
\section{Related Work}
\label{sec:related}

\subsection{DLO state estimation without vision}
\begin{outline}
  \item Vision is the dominant modality for DLO perception and
    manipulation~\cite{sanchez2018survey}, but it fails under the occlusions we
    target. Force/torque (F/T) and tactile sensing are therefore common in DLO
    insertion and
    routing~\cite{kr-aecsagphi-22,nkot-sditfw-97,slk-ftmdloeffi-22,wjly-cratmp-22,zdpp-dhattfmn-19,lvw-crcrisb-24}.
  \item For an elastic rod in static equilibrium, the shape is determined by
    the wrench at one end~\cite{bretl2014quasistatic}; this enabled real-time,
    vision-free shape estimation of a wire held by two arms from F/T
    sensing~\cite{mishani2022realtime}.
  \item \citet{slk-ftmdloeffi-22} manipulate a DLO by keeping it under tension,
    assuming it is straight when tensioned and that gravity makes free segments
    vertical; F/T readings detect the tension and feed an analytic model for
    estimating the DLO state.
  \item Gap: these works assume a contact-free rod held at both ends, or a DLO
    kept taut. An environment contact that bends the rod adds an unknown
    reaction force that breaks the wrench-to-shape mapping; we add the contact
    location to recover it.
\end{outline}

\subsection{DLO insertion and threading}
\begin{outline}
  \item \citet{nkot-sditfw-97} insert a flexible wire using a stereo camera and
    an F/T sensor, first straightening its free end; the insertion strategy is
    standard, and deformation caused by contact with the environment is not
    discussed.
  \item Tight-tolerance threading of string and rope through
    openings~\cite{wang2015online}; learned DLO insertion with flexibility
    estimated from visual cues~\cite{lc-ldloilfevc-24}.
  \item Tactile insertion: \citet{zdpp-dhattfmn-19} learn to guide a deformed
    DLO through a hole using a 16-taxel tactile sensor, without assuming
    compliance or deforming the DLO. \citet{wjly-cratmp-22} build a library of
    tactile-driven primitives (cable following, pivoting, weaving, insertion)
    using a GelSight sensor, designed for the NIST Assembly Task Board; each
    primitive is a state machine whose transitions are triggered by tactile
    input, and insertion is standard: push until the force drops and a rise in
    torque shows that rotation is impossible. Tactile feedback also enables
    following a cable in the grip~\cite{she2021cable}.
  \item \citet{lvw-crcrisb-24} use reinforcement learning (RL) with F/T sensing
    to push compliant objects, including DLOs, into tunnels of unknown shape; a
    Dynamic Safety Lock computes the forces acting on the end effector to allow
    some pushing but stops when the resistance is too high. This is the only
    work we know of with non-straight passages, but its bends do not exceed
    roughly $\pi/6$, and it addresses the uncertainty of the passage shape
    rather than motion planning or control.
  \item Gap: insertion either relies on vision or pushes along a straight or
    mildly bent passage until the force changes; contact with the passage is
    treated as resistance or uncertainty, not used to redirect the tip, so
    angular threading behind an occluding edge is not addressed.
\end{outline}

\subsection{Exploiting environment contacts}
\begin{outline}
  \item Extrinsic dexterity uses external contacts to extend what a simple
    gripper can do~\cite{dafle2014extrinsic}. For peg-in-hole insertion,
    \citet{kr-aecsagphi-22} estimate the contact line between the peg and the
    hole with a GelSight sensor and feed it to an RL model that selects the
    next action.
  \item For DLOs, environmental contacts and fixtures have been used to shape
    cables~\cite{zhu2020environmental} and to manipulate them blindly with
    force feedback~\cite{slk-ftmdloeffi-22}.
  \item Contacts along slender elastic bodies can be detected and localized
    from kinematics, as done for continuum robots~\cite{bajo2012kinematics}.
  \item Gap: contacts are used to shape or hold the DLO body, or are sensed on a
    rigid peg, not used as a joint for steering a \emph{free tip} to a target
    pose with haptic feedback.
    \todo{Confirm this novelty claim against recent work (2024--2026).}
\end{outline}

%% ============================================================================
%% III. Problem Statement                                     target: ~0.5 page
%% Stage 1: the problem statement must be fully developed (a concrete technical
%% formulation, not a high-level argument).
%% ============================================================================
\section{Problem Statement}
\label{sec:problem}

\begin{figure}[t]
  \centering
  \figplaceholder{\linewidth}{3.6cm}{%
    \textbf{Fig.~2 (notation).} Gripper pose $\gpose$, rod of length $L$,
    contact at arc length $s_{\mathrm{c}}$ and ring angle $\cang$, wrist wrench
    $\wrench$, tip pose $\tpose$, and target $\tposegoal$.}
  \caption{Problem setup and notation.}
  \label{fig:notation}
\end{figure}

\begin{outline}
  \item \textbf{Object.} A naturally straight, superelastic Nitinol rod of
    length $L$, modeled as an inextensible Kirchhoff rod under quasi-static
    manipulation~\cite{bretl2014quasistatic}, grasped rigidly at one end; the
    free tip carries a tool. \todo{Rod length, diameter, material constants.}
  \item \textbf{Robot and sensing.} Gripper pose $\gpose\in\SE{3}$ ($\SE{2}$
    in planar experiments) from forward kinematics; wrist wrench
    $\wrench=(\vect{f},\vect{\tau})\in\R^{6}$ from a force/torque sensor; a
    sensorized environment edge (ring) that reports contact and the contact
    angle $\cang$ along the ring (\cref{fig:notation}).
  \item \textbf{Contact.} The rod touches the edge at a single point at an
    unknown arc length $s_{\mathrm{c}}\in(0,L)$; the edge applies an unknown
    reaction force (normal and friction). \todo{Sticking vs.\ sliding contact;
    friction model.}
  \item \textbf{Observation} at time $t$: $\vect{o}_t=(\gpose,\wrench,\cpt)$,
    where $\cpt$ holds the contact state and $\cang$. No vision is used during
    execution; an RGB-D camera provides ground truth only.
  \item \textbf{Reachability.} Let $\Rfree$ and $\Rcon$ be the sets of tip
    poses reachable in the workspace $\mathcal{W}$ without and with an edge
    contact, respectively.
  \item \textbf{Assumptions.} Quasi-static motion, a single contact, known edge
    geometry, deformation within the rod's elastic range.
    \todo{Gravity: modeled or negligible?}
\end{outline}

\begin{problem}[Haptic tip-pose estimation]\label{prob:estimation}
Given the observation $\vect{o}_t$, estimate the tip pose, $\tposehat\approx\tpose$.
\end{problem}

\begin{problem}[Contact-aided tip manipulation]\label{prob:manipulation}
Given a target $\tposegoal\in\Rcon\setminus\Rfree$, find gripper motions that
bring $\tpose$ to within position and orientation tolerances
$(\epsilon_p,\epsilon_\theta)$ of $\tposegoal$ by using a contact with the
edge, while avoiding other collisions and keeping the rod within its elastic
range.
\end{problem}

%% ============================================================================
%% IV. Hypotheses                                             target: ~0.4 page
%% Stage-1 reviewers score (i) the value of the hypotheses and (ii) whether the
%% planned evaluation can support OR falsify them. Make each hypothesis
%% testable and state what result would refute it. Empirical hypotheses are OK.
%% Reference them as \ref{hyp:...} (prints H1, H2, ...).
%% ============================================================================
\section{Hypotheses}
\label{sec:hypotheses}

\todo{Refine the wording and add quantitative thresholds once the setup is fixed.}

\begin{hypothesis}[Reachability]\label{hyp:reach}
Allowing a single contact between the rod and an environment edge strictly
enlarges the set of reachable tip poses, in particular the range of tip
directions, compared with contact-free manipulation in the same constrained
workspace ($\Rfree\subsetneq\Rcon$).
\end{hypothesis}
\begin{outline}
  \item Rationale: the contact acts as an ad-hoc joint (pivot), and its
    location along the rod is controlled by feeding the rod over the edge.
  \item Refuted if, for the tested clearances, the tip directions reachable
    with contact do not exceed those reachable without it (simulation and
    hardware).
\end{outline}

\begin{hypothesis}[Vision-free estimation]\label{hyp:estimate}
Gripper pose, wrist wrench, and contact location suffice to estimate the tip
pose without vision, and the contact location significantly reduces the error
compared with gripper pose and wrench alone.
\end{hypothesis}
\begin{outline}
  \item Rationale: without contact, the wrench at the gripped end determines
    the rod's shape~\cite{bretl2014quasistatic,mishani2022realtime}; a contact
    adds an unknown reaction force, and knowing where it acts recovers the
    missing information.
  \item Refuted if the tip error on unseen configurations exceeds the task
    tolerance, or if adding the contact location gives no significant
    improvement (input ablation, \cref{sec:experiments}).
\end{outline}

\begin{hypothesis}[Contact-aided steering]\label{hyp:steer}
Closed-loop steering with the haptic tip-pose estimate brings the tip to
target poses behind the edge with a higher success rate than contact-free and
open-loop strategies.
\end{hypothesis}
\begin{outline}
  \item Refuted if the success rate and final error are not better than those
    of the baselines on the same targets.
\end{outline}

\begin{outline}
  \item \textbf{Expected insights:} when contact helps and by how much (as a
    function of rod stiffness, clearance, and contact location), and which
    haptic signals carry the tip-pose information.
\end{outline}

%% ============================================================================
%% V. Methodology                                             target: ~1.3 pages
%% Stage 1: the methodology must be fully developed, including the specific
%% design choices in the system architecture and the experiments.
%% ============================================================================
\section{Methodology}
\label{sec:method}

\begin{figure}[t]
  \centering
  \figplaceholder{\linewidth}{3.2cm}{%
    \textbf{Fig.~3 (pipeline).} Inputs $\gpose$, $\wrench$, $\cpt$
    $\rightarrow$ tip-pose estimator $\rightarrow$ $\tposehat$ $\rightarrow$
    contact-aware steering $\rightarrow$ gripper motion command.}
  \caption{Overview of the approach.}
  \label{fig:pipeline}
\end{figure}

\subsection{Contact as an ad-hoc joint}
\label{sec:method:reach}
\begin{outline}
  \item Model the rod as a quasi-static Kirchhoff rod~\cite{bretl2014quasistatic};
    the edge contact is a point constraint that transmits force but no moment,
    \ie a passive joint whose location along the rod, $s_{\mathrm{c}}$, changes
    as the rod is fed over the edge.
  \item Compute $\Rfree$ and $\Rcon$ by sampling gripper poses in simulation
    (\eg with discrete elastic rods~\cite{bergou2008der}) and validate on
    hardware; this tests \ref{hyp:reach}.
  \item \todo{Analytic result for the planar case (elastica)? Bound on the
    achievable change in tip direction as a function of clearance.}
\end{outline}

\subsection{Sensorized edge}
\label{sec:method:edge}
\begin{outline}
  \item Conductive TPU ring around the opening that reports contact and the
    angular contact location $\cang$ along the ring.
    \todo{Sensing principle, resolution, latency, calibration.}
  \item Design choice: measure the contact location directly instead of
    inferring it from kinematics or wrench (\cf~\cite{bajo2012kinematics}), to
    isolate the value of contact information (\ref{hyp:estimate}).
\end{outline}

\subsection{Haptic tip-pose estimation}
\label{sec:method:estimation}
\begin{outline}
  \item Learn $f_{\theta}:\vect{o}_t\mapsto\tposehat$ from data collected on
    the real system (\cref{sec:experiments}); camera ground truth is used only
    as training labels.
  \item Inputs: $\gpose$, the tared and gravity-compensated wrench $\wrench$,
    the contact state, and $\cang$. Outputs: tip position and direction
    ($(x,y,\theta)$ in the planar case).
  \item Candidate models to compare: MLP; a temporal model over a short history
    (\eg 1D-CNN or GRU) to capture friction and hysteresis; a physics-informed
    model that learns a residual on top of a rod-model
    prediction~\cite{mishani2022realtime}. \todo{Final list of models.}
  \item Loss: weighted position and orientation error; uncertainty from an
    ensemble to gate the controller. \todo{Decide.}
\end{outline}

\subsection{Contact-aware steering}
\label{sec:method:steering}
\begin{outline}
  \item Phases: (i) approach and make contact at a chosen edge location;
    (ii) pivot about the contact to rotate the tip direction; (iii) feed the
    rod to translate the tip to $\tposegoal$.
  \item Closed loop on $\tposehat$: local Jacobian-based servoing (Jacobian
    estimated online or from the learned model), with limits on contact force
    and rod curvature. \todo{Alternative: sampling-based planning in simulation
    with feedback correction.}
\end{outline}

\subsection{Implementation}
\label{sec:method:implementation}
\begin{outline}
  \item ROS~1 (Noetic) with MoveIt for motion planning; synchronized logging of
    joint states, wrench, contact signal, and camera frames.
    \todo{Synchronization method and rates.}
\end{outline}

%% ============================================================================
%% VI. Experimental Plan and Preliminary Results              target: ~1.3 pages
%% Stage 1: experiments may be missing or preliminary, but the plan must show
%% how each hypothesis will be supported or falsified (reviewers score the
%% "feasibility of conclusion"). Stage 2: rename to "Experiments" and report
%% the experiments required by the Extension Charter.
%% Anonymity: no lab names or logos in photos or text.
%% ============================================================================
\section{Experimental Plan and Preliminary Results}
\label{sec:experiments}

\begin{figure}[t]
  \centering
  \figplaceholder{\linewidth}{4cm}{%
    \textbf{Fig.~4 (setup).} Robot arm with wrist force/torque sensor and
    gripper holding the Nitinol rod; sensorized ring (edge); RGB-D camera used
    for ground truth only.}
  \caption{Experimental setup.}
  \label{fig:setup}
\end{figure}

\subsection{Setup}
\label{sec:exp:setup}
\begin{outline}
  \item Yaskawa Motoman SDA10F (one arm), Bota SenseONE EtherCAT force/torque
    sensor at the wrist, parallel gripper \todo{model}, Nitinol rod
    \todo{length, diameter}, conductive-TPU sensorized ring, Intel RealSense
    D435 RGB-D camera for ground truth only \todo{tip marker / tracking method}.
  \item Planar (2D) experiments first, then spatial (3D).
\end{outline}

\subsection{Data collection and protocol}
\label{sec:exp:protocol}
\begin{outline}
  \item Random and structured gripper trajectories that cover contact and
    no-contact states over a range of deformations; each sample is
    $(\gpose,\wrench,\cpt,\tpose)$.
  \item Train/test split by trajectory (testing on unseen configurations);
    repeated trials to quantify sensor noise and hysteresis.
  \item Mitigations for sensor noise, inaccurate contact information, and
    synchronization errors: force/torque tare and calibration, time
    synchronization, camera--robot extrinsic calibration, repeated controlled
    trials.
  \item \todo{Dataset size, number of rods and edges, range of contact
    locations.}
\end{outline}

\subsection{Metrics and baselines}
\label{sec:exp:metrics}
\begin{outline}
  \item Estimation: tip position error [mm] and direction error [deg] on
    held-out data.
  \item Manipulation: success rate within $(\epsilon_p,\epsilon_\theta)$, final
    error, time to target, peak contact force.
  \item Baselines and ablations: inputs ($\gpose$; $\gpose{+}\wrench$;
    $\gpose{+}\wrench{+}\cpt$); model comparison; a rod-model (physics)
    estimator; contact-free and open-loop steering.
\end{outline}

\begin{table}[t]
  \centering
  \caption{How each hypothesis is tested}
  \label{tab:plan}
  \footnotesize
  \begin{tabular}{@{}p{0.06\linewidth}p{0.40\linewidth}p{0.44\linewidth}@{}}
    \toprule
     & Experiment & Supported if / refuted if \\
    \midrule
    \ref{hyp:reach}    & Reachable tip directions with vs.\ without contact
                         (simulation and hardware)
                       & Larger range with contact / no gain \\
    \ref{hyp:estimate} & Estimation error on unseen data; input ablation
                       & Error within tolerance and significant gain from
                         contact input / not \\
    \ref{hyp:steer}    & Steering to targets behind the edge vs.\ baselines
                       & Higher success, lower final error / not \\
    \bottomrule
  \end{tabular}
\end{table}

\subsection{Preliminary results}
\label{sec:exp:prelim}

\begin{table}[t]
  \centering
  \caption{Preliminary tip-pose estimation error (planar, held-out data)}
  \label{tab:prelim}
  \footnotesize
  \begin{tabular}{@{}lcc@{}}
    \toprule
    Inputs & Position [mm] & Direction [deg] \\
    \midrule
    $\gpose$                    & -- & -- \\
    $\gpose+\wrench$            & -- & -- \\
    $\gpose+\wrench+\cpt$       & -- & -- \\
    \bottomrule
  \end{tabular}
\end{table}

\begin{outline}
  \item \todo{First planar dataset: size, error statistics (\cref{tab:prelim}),
    and what they suggest about \ref{hyp:estimate}.}
  \item \todo{Reachability in simulation for \ref{hyp:reach}.}
\end{outline}

\subsection{Planned experiments (Stage 2)}
\label{sec:exp:planned}
\begin{outline}
  \item Spatial (3D) experiments; closed-loop steering trials
    (\ref{hyp:steer}); generalization to other rods and edge geometries.
\end{outline}

%% ============================================================================
%% VII. Conclusion and Expected Insights                     target: ~0.25 page
%% Stage 1 asks for the expected scientific insights, not only results.
%% ============================================================================
\section{Conclusion and Expected Insights}
\label{sec:conclusion}

\begin{outline}
  \item Summary: a contact point with an environment edge, used as an ad-hoc
    joint, together with vision-free tip-pose estimation, for steering the tip
    of an elastic rod in constrained, occluded spaces.
  \item Expected insights: the conditions under which contact enlarges
    reachability; how much tip-pose information the wrench and the contact
    location carry; failure modes (sliding, loss of contact, multiple
    contacts).
  \item Limitations: a single contact, quasi-static motion, known edge
    geometry, one rod material; extension to cables with low bending stiffness
    is open.
  \item Stage-2 plan: 3D experiments and closed-loop steering trials, to be
    refined by the Extension Charter.
\end{outline}


%% Acknowledgments are not allowed in the anonymous submission.
\ifanonymous\else
\section*{Acknowledgments}
\todo{Funding and acknowledgments (camera-ready only).}
\fi

%% Use plainnat to work nicely with natbib (RSS template).
\bibliographystyle{plainnat}
\bibliography{references}

\end{document}
"
%% ---------------------------------------------------------------------------
\providecommand*{\papermode}{draft}
\newcommand*{\paperid}{[add your ID here]} % OpenReview paper ID (Stage 1)
\newcommand*{\papertitle}{Elastic Rod Manipulation Using a Contact Point}

%% ---------------------------------------------------------------------------
%% Packages from the RSS template (hyperref is loaded further down, after the
%% packages it must follow)
%% ---------------------------------------------------------------------------
\usepackage{times}
% numbers option provides compact numerical references in the text.
\usepackage[numbers]{natbib}
\usepackage{multicol}

%% ---------------------------------------------------------------------------
%% Additional packages (none of them changes fonts, margins, or spacing)
%% Avoid: subcaption/caption, enumitem, geometry, mathptmx/newtx, titlesec --
%% they alter the IEEEtran layout.
%% ---------------------------------------------------------------------------
\usepackage{amsmath}
\usepackage{amssymb}
\usepackage{bm}
\usepackage{graphicx}
\graphicspath{{figures/}}
\usepackage[caption=false,font=footnotesize]{subfig} % IEEEtran-safe subfigures
\usepackage{booktabs}
\usepackage{multirow}
\usepackage{siunitx}  % \SI{2}{\milli\metre}, \SI{5}{\degree}
\usepackage{xcolor}
\usepackage{algorithm}
\usepackage{algpseudocode}
\usepackage{xspace}
\usepackage{etoolbox}

\usepackage[bookmarks=true]{hyperref} % RSS template package
\usepackage[capitalise]{cleveref}     % must come after hyperref
\crefname{section}{Sec.}{Secs.}
\crefname{subsection}{Sec.}{Secs.}
\crefname{figure}{Fig.}{Figs.}
\crefname{table}{Table}{Tables}
\crefname{equation}{Eq.}{Eqs.}
\crefname{algorithm}{Alg.}{Algs.}

%% ---------------------------------------------------------------------------
%% Build-mode switches
%% ---------------------------------------------------------------------------
\newif\ifanonymous
\newif\ifshownotes
\ifdefstring{\papermode}{submission}{\anonymoustrue\shownotesfalse}{%
  \ifdefstring{\papermode}{camera}{\anonymousfalse\shownotesfalse}{%
    \anonymousfalse\shownotestrue}}% draft (also the fallback)
\typeout{*** Paper mode: \papermode}

%% PDF metadata (must be anonymous in the review stages)
\hypersetup{
  pdftitle    = {\papertitle},
  pdfkeywords = {deformable linear objects; elastic rods; contact; haptic sensing; manipulation},
}
\ifanonymous
  \hypersetup{pdfauthor={}}
\else
  \hypersetup{pdfauthor={Stav Ashur, Yuval Cassuto, Avishai Sintov}}
\fi

%% ---------------------------------------------------------------------------
%% Draft notes: visible only in draft mode, removed in submission/camera modes
%% (do not also load the todonotes package: it defines \todo as well)
%% ---------------------------------------------------------------------------
\definecolor{notecolor}{RGB}{200,0,0}
\definecolor{outlinecolor}{RGB}{0,70,150}
\DeclareRobustCommand{\draftnote}[3]{%
  \ifshownotes\textcolor{#1}{[\textbf{#2:} #3]}\else\ifhmode\unskip\fi\fi}
\newcommand{\todo}[1]{\draftnote{notecolor}{TODO}{#1}}
\newcommand{\stav}[1]{\draftnote{orange}{Stav}{#1}}
\newcommand{\yuval}[1]{\draftnote{teal}{Yuval}{#1}}
\newcommand{\avishai}[1]{\draftnote{violet}{Avishai}{#1}}

%% Skeleton outlines: bullet lists of planned content. Replace each outline
%% block with prose as you write; they stay blue in every mode so leftovers
%% are easy to spot, and non-draft builds log how many are left.
\newcounter{outlineblocks}
\newenvironment{outline}{%
  \stepcounter{outlineblocks}\par\color{outlinecolor}\begin{itemize}}%
  {\end{itemize}}
\AtEndDocument{\ifshownotes\else\ifnum\value{outlineblocks}>0
  \typeout{WARNING: \arabic{outlineblocks} skeleton outline block(s) left in the paper}%
\fi\fi}

%% \figplaceholder{<width>}{<height>}{<description>}: framed stand-in for a
%% figure that does not exist yet. Replace with \includegraphics{...}.
\newcommand{\figplaceholder}[3]{%
  \fbox{\parbox[c][#2][c]{\dimexpr #1-2\fboxsep-2\fboxrule\relax}{%
    \centering\footnotesize\color{outlinecolor}#3}}}

%% ---------------------------------------------------------------------------
%% Hypotheses and problem statements
%%   \begin{hypothesis}[Short name]\label{hyp:x} ... \end{hypothesis}
%%   \ref{hyp:x} prints "H1", "H2", ...
%% ---------------------------------------------------------------------------
\newcounter{hypothesis}
\renewcommand{\thehypothesis}{H\arabic{hypothesis}}
\crefname{hypothesis}{hypothesis}{hypotheses}
\newenvironment{hypothesis}[1][]{%
  \par\smallskip\refstepcounter{hypothesis}\noindent
  \textbf{\thehypothesis\ifstrempty{#1}{}{ (#1)}:}\hspace{0.5em}\itshape\ignorespaces}%
  {\par\smallskip}
\newtheorem{problem}{Problem}
\crefname{problem}{Problem}{Problems}

%% ---------------------------------------------------------------------------
%% Notation
%% ---------------------------------------------------------------------------
\newcommand{\R}{\mathbb{R}}
\newcommand{\SE}[1]{\mathrm{SE}(#1)}
\newcommand{\SO}[1]{\mathrm{SO}(#1)}
\newcommand{\vect}[1]{\bm{#1}}
\newcommand{\gpose}{\vect{T}_{\mathrm{g}}}          % gripper (end-effector) pose
\newcommand{\tpose}{\vect{T}_{\mathrm{tip}}}        % DLO tip pose
\newcommand{\tposehat}{\hat{\vect{T}}_{\mathrm{tip}}} % estimated tip pose
\newcommand{\tposegoal}{\vect{T}^{*}_{\mathrm{tip}}} % target tip pose
\newcommand{\wrench}{\vect{w}}                      % wrist F/T reading (f, tau)
\newcommand{\cpt}{\vect{c}}                         % contact measurement
\newcommand{\cang}{\phi_{\mathrm{c}}}               % contact angle on the ring
\newcommand{\Rfree}{\mathcal{R}_{\mathrm{free}}}    % reachable set, no contact
\newcommand{\Rcon}{\mathcal{R}_{\mathrm{c}}}        % reachable set, with contact

\newcommand{\eg}{e.g.,\xspace}
\newcommand{\ie}{i.e.,\xspace}
\newcommand{\etal}{et~al.\xspace}
\newcommand{\wrt}{w.r.t.\xspace}
\newcommand{\cf}{cf.\xspace}

\begin{document}

\title{\papertitle}

\ifanonymous
  % Author block required by the RSS template for anonymous review
  \author{Author Names Omitted for Anonymous Review. Paper-ID \paperid}
\else
  \author{\authorblockN{Stav Ashur, Yuval Cassuto, and Avishai Sintov}
  \authorblockA{Tel Aviv University, Tel Aviv, Israel}}
  % TODO (camera-ready): departments and e-mail addresses
\fi

\maketitle

\begin{abstract}
Sensors and devices such as cameras and brushes can be mounted on the tips of
elastic rods and other deformable linear objects (DLOs), allowing them to be
inserted into and manipulated within constrained spaces. Using the right
contact point enables the DLO to reach poses that are infeasible when no
contact with the environment is acceptable. Such a contact point serves as an
ad-hoc joint, facilitating direction changes of the DLO's tip that are
impossible if no collisions are tolerated. In this work, we propose the use of
haptic sensors to perform manipulation that brings the tip of the DLO to a
target pose in a constrained environment by utilizing a contact point of the
DLO with an environment edge.
\todo{Add 1--2 sentences on the hypotheses and methodology (vision-free tip-pose
estimation from gripper pose, wrist force/torque, and contact location;
contact-aware steering) and on the preliminary results.}
\end{abstract}

\IEEEpeerreviewmaketitle

%% Sections, in order
%% ============================================================================
%% I. Introduction                               target: ~0.75 page incl. Fig. 1
%% Motivate the problem, state the gap, preview the hypotheses (Sec. IV) and
%% the methodology (Sec. V), and list the contributions.
%% ============================================================================
\section{Introduction}
\label{sec:intro}

\begin{outline}
  \item \textbf{Motivation.} Cameras, brushes, and other tools are mounted on
    the tips of elastic rods and other DLOs and inserted into constrained
    spaces, \eg for inspection and repair inside engines~\cite{dong2017slender}.
    \todo{1--2 more application examples.}
  \item \textbf{Challenge.} The rod's shape depends on mechanics that are hard
    to model~\cite{sanchez2018survey}, and once inserted, the rod is hidden by
    the environment or by the robot itself, so visual state estimation is
    unreliable.
  \item \textbf{Key idea.} Contact with the environment is usually avoided.
    Instead, we use a contact point between the rod and an environment edge as
    an \emph{ad-hoc joint}: pivoting about it changes the tip direction in ways
    that are impossible if no collisions are tolerated (\cref{fig:teaser}).
  \item \textbf{Problem.} Bring the tip of a gripped elastic rod to a target
    pose in a constrained environment by exploiting a contact with an
    environment edge, using only haptic sensing (robot proprioception, wrist
    force/torque, contact location) during execution (\cref{sec:problem}).
  \item \textbf{Approach.} A sensorized edge reports where the rod touches it;
    a learned estimator maps gripper pose, wrench, and contact location to the
    tip pose; a contact-aware controller steers the tip to the target
    (\cref{sec:method}).
\end{outline}

% Fig. 1 is defined here, not at the top of the section, so that it lands at
% the top of the right column of page 1. LaTeX puts a [t] float at the top of
% the left column (above the abstract) if that column still has room when the
% float is reached; if this happens after edits, move this block further down.
\begin{figure}[t]
  \centering
  \figplaceholder{\linewidth}{4.2cm}{%
    \textbf{Fig.~1 (teaser).} A robot grasps an elastic rod whose tip carries a
    tool. The rod rests on an environment edge (sensorized ring) that acts as an
    ad-hoc joint, letting the tip turn toward a target pose that is unreachable
    without contact. Inset: reachable tip directions with vs.\ without contact.}
  \caption{Using a contact point with an environment edge as an ad-hoc joint
    lets the tip of a gripped elastic rod reach poses that are infeasible
    without contact.}
  \label{fig:teaser}
\end{figure}

\begin{outline}
  \item \textbf{Hypotheses} (\cref{sec:hypotheses}): contact enlarges the set
    of reachable tip poses (\ref{hyp:reach}); haptic signals suffice to estimate
    the tip pose without vision (\ref{hyp:estimate}); closed-loop haptic
    steering reaches targets that contact-free or open-loop strategies miss
    (\ref{hyp:steer}).
  \item \textbf{Contributions}:
    \begin{itemize}
      \item a formulation of DLO tip manipulation that uses a single
        environment contact as an ad-hoc joint, with a reachability analysis;
      \item a vision-free tip-pose estimator that uses proprioception,
        force/torque, and contact location;
      \item an experimental methodology and dataset with a Nitinol rod, from
        planar (2D) to spatial (3D) settings;
      \item preliminary results on the hypotheses.
        \todo{Update after the experiments.}
    \end{itemize}
\end{outline}

%% ============================================================================
%% II. Related Work                                           target: ~0.5 page
%% Stage 1: cite only the work most directly tied to the central problem (what
%% is needed to establish novelty). Stage 2 requires a complete section.
%% Anonymity: cite the authors' own prior work in the third person.
%% RSS 2027 checks references automatically and desk-rejects hallucinated
%% citations, so check every entry you add to references.bib.
%% Merged with the notes in material/2_Related_work.tex (earlier document).
%% ============================================================================
\section{Related Work}
\label{sec:related}

\subsection{DLO state estimation without vision}
\begin{outline}
  \item Vision is the dominant modality for DLO perception and
    manipulation~\cite{sanchez2018survey}, but it fails under the occlusions we
    target. Force/torque (F/T) and tactile sensing are therefore common in DLO
    insertion and
    routing~\cite{kr-aecsagphi-22,nkot-sditfw-97,slk-ftmdloeffi-22,wjly-cratmp-22,zdpp-dhattfmn-19,lvw-crcrisb-24}.
  \item For an elastic rod in static equilibrium, the shape is determined by
    the wrench at one end~\cite{bretl2014quasistatic}; this enabled real-time,
    vision-free shape estimation of a wire held by two arms from F/T
    sensing~\cite{mishani2022realtime}.
  \item \citet{slk-ftmdloeffi-22} manipulate a DLO by keeping it under tension,
    assuming it is straight when tensioned and that gravity makes free segments
    vertical; F/T readings detect the tension and feed an analytic model for
    estimating the DLO state.
  \item Gap: these works assume a contact-free rod held at both ends, or a DLO
    kept taut. An environment contact that bends the rod adds an unknown
    reaction force that breaks the wrench-to-shape mapping; we add the contact
    location to recover it.
\end{outline}

\subsection{DLO insertion and threading}
\begin{outline}
  \item \citet{nkot-sditfw-97} insert a flexible wire using a stereo camera and
    an F/T sensor, first straightening its free end; the insertion strategy is
    standard, and deformation caused by contact with the environment is not
    discussed.
  \item Tight-tolerance threading of string and rope through
    openings~\cite{wang2015online}; learned DLO insertion with flexibility
    estimated from visual cues~\cite{lc-ldloilfevc-24}.
  \item Tactile insertion: \citet{zdpp-dhattfmn-19} learn to guide a deformed
    DLO through a hole using a 16-taxel tactile sensor, without assuming
    compliance or deforming the DLO. \citet{wjly-cratmp-22} build a library of
    tactile-driven primitives (cable following, pivoting, weaving, insertion)
    using a GelSight sensor, designed for the NIST Assembly Task Board; each
    primitive is a state machine whose transitions are triggered by tactile
    input, and insertion is standard: push until the force drops and a rise in
    torque shows that rotation is impossible. Tactile feedback also enables
    following a cable in the grip~\cite{she2021cable}.
  \item \citet{lvw-crcrisb-24} use reinforcement learning (RL) with F/T sensing
    to push compliant objects, including DLOs, into tunnels of unknown shape; a
    Dynamic Safety Lock computes the forces acting on the end effector to allow
    some pushing but stops when the resistance is too high. This is the only
    work we know of with non-straight passages, but its bends do not exceed
    roughly $\pi/6$, and it addresses the uncertainty of the passage shape
    rather than motion planning or control.
  \item Gap: insertion either relies on vision or pushes along a straight or
    mildly bent passage until the force changes; contact with the passage is
    treated as resistance or uncertainty, not used to redirect the tip, so
    angular threading behind an occluding edge is not addressed.
\end{outline}

\subsection{Exploiting environment contacts}
\begin{outline}
  \item Extrinsic dexterity uses external contacts to extend what a simple
    gripper can do~\cite{dafle2014extrinsic}. For peg-in-hole insertion,
    \citet{kr-aecsagphi-22} estimate the contact line between the peg and the
    hole with a GelSight sensor and feed it to an RL model that selects the
    next action.
  \item For DLOs, environmental contacts and fixtures have been used to shape
    cables~\cite{zhu2020environmental} and to manipulate them blindly with
    force feedback~\cite{slk-ftmdloeffi-22}.
  \item Contacts along slender elastic bodies can be detected and localized
    from kinematics, as done for continuum robots~\cite{bajo2012kinematics}.
  \item Gap: contacts are used to shape or hold the DLO body, or are sensed on a
    rigid peg, not used as a joint for steering a \emph{free tip} to a target
    pose with haptic feedback.
    \todo{Confirm this novelty claim against recent work (2024--2026).}
\end{outline}

%% ============================================================================
%% III. Problem Statement                                     target: ~0.5 page
%% Stage 1: the problem statement must be fully developed (a concrete technical
%% formulation, not a high-level argument).
%% ============================================================================
\section{Problem Statement}
\label{sec:problem}

\begin{figure}[t]
  \centering
  \figplaceholder{\linewidth}{3.6cm}{%
    \textbf{Fig.~2 (notation).} Gripper pose $\gpose$, rod of length $L$,
    contact at arc length $s_{\mathrm{c}}$ and ring angle $\cang$, wrist wrench
    $\wrench$, tip pose $\tpose$, and target $\tposegoal$.}
  \caption{Problem setup and notation.}
  \label{fig:notation}
\end{figure}

\begin{outline}
  \item \textbf{Object.} A naturally straight, superelastic Nitinol rod of
    length $L$, modeled as an inextensible Kirchhoff rod under quasi-static
    manipulation~\cite{bretl2014quasistatic}, grasped rigidly at one end; the
    free tip carries a tool. \todo{Rod length, diameter, material constants.}
  \item \textbf{Robot and sensing.} Gripper pose $\gpose\in\SE{3}$ ($\SE{2}$
    in planar experiments) from forward kinematics; wrist wrench
    $\wrench=(\vect{f},\vect{\tau})\in\R^{6}$ from a force/torque sensor; a
    sensorized environment edge (ring) that reports contact and the contact
    angle $\cang$ along the ring (\cref{fig:notation}).
  \item \textbf{Contact.} The rod touches the edge at a single point at an
    unknown arc length $s_{\mathrm{c}}\in(0,L)$; the edge applies an unknown
    reaction force (normal and friction). \todo{Sticking vs.\ sliding contact;
    friction model.}
  \item \textbf{Observation} at time $t$: $\vect{o}_t=(\gpose,\wrench,\cpt)$,
    where $\cpt$ holds the contact state and $\cang$. No vision is used during
    execution; an RGB-D camera provides ground truth only.
  \item \textbf{Reachability.} Let $\Rfree$ and $\Rcon$ be the sets of tip
    poses reachable in the workspace $\mathcal{W}$ without and with an edge
    contact, respectively.
  \item \textbf{Assumptions.} Quasi-static motion, a single contact, known edge
    geometry, deformation within the rod's elastic range.
    \todo{Gravity: modeled or negligible?}
\end{outline}

\begin{problem}[Haptic tip-pose estimation]\label{prob:estimation}
Given the observation $\vect{o}_t$, estimate the tip pose, $\tposehat\approx\tpose$.
\end{problem}

\begin{problem}[Contact-aided tip manipulation]\label{prob:manipulation}
Given a target $\tposegoal\in\Rcon\setminus\Rfree$, find gripper motions that
bring $\tpose$ to within position and orientation tolerances
$(\epsilon_p,\epsilon_\theta)$ of $\tposegoal$ by using a contact with the
edge, while avoiding other collisions and keeping the rod within its elastic
range.
\end{problem}

%% ============================================================================
%% IV. Hypotheses                                             target: ~0.4 page
%% Stage-1 reviewers score (i) the value of the hypotheses and (ii) whether the
%% planned evaluation can support OR falsify them. Make each hypothesis
%% testable and state what result would refute it. Empirical hypotheses are OK.
%% Reference them as \ref{hyp:...} (prints H1, H2, ...).
%% ============================================================================
\section{Hypotheses}
\label{sec:hypotheses}

\todo{Refine the wording and add quantitative thresholds once the setup is fixed.}

\begin{hypothesis}[Reachability]\label{hyp:reach}
Allowing a single contact between the rod and an environment edge strictly
enlarges the set of reachable tip poses, in particular the range of tip
directions, compared with contact-free manipulation in the same constrained
workspace ($\Rfree\subsetneq\Rcon$).
\end{hypothesis}
\begin{outline}
  \item Rationale: the contact acts as an ad-hoc joint (pivot), and its
    location along the rod is controlled by feeding the rod over the edge.
  \item Refuted if, for the tested clearances, the tip directions reachable
    with contact do not exceed those reachable without it (simulation and
    hardware).
\end{outline}

\begin{hypothesis}[Vision-free estimation]\label{hyp:estimate}
Gripper pose, wrist wrench, and contact location suffice to estimate the tip
pose without vision, and the contact location significantly reduces the error
compared with gripper pose and wrench alone.
\end{hypothesis}
\begin{outline}
  \item Rationale: without contact, the wrench at the gripped end determines
    the rod's shape~\cite{bretl2014quasistatic,mishani2022realtime}; a contact
    adds an unknown reaction force, and knowing where it acts recovers the
    missing information.
  \item Refuted if the tip error on unseen configurations exceeds the task
    tolerance, or if adding the contact location gives no significant
    improvement (input ablation, \cref{sec:experiments}).
\end{outline}

\begin{hypothesis}[Contact-aided steering]\label{hyp:steer}
Closed-loop steering with the haptic tip-pose estimate brings the tip to
target poses behind the edge with a higher success rate than contact-free and
open-loop strategies.
\end{hypothesis}
\begin{outline}
  \item Refuted if the success rate and final error are not better than those
    of the baselines on the same targets.
\end{outline}

\begin{outline}
  \item \textbf{Expected insights:} when contact helps and by how much (as a
    function of rod stiffness, clearance, and contact location), and which
    haptic signals carry the tip-pose information.
\end{outline}

%% ============================================================================
%% V. Methodology                                             target: ~1.3 pages
%% Stage 1: the methodology must be fully developed, including the specific
%% design choices in the system architecture and the experiments.
%% ============================================================================
\section{Methodology}
\label{sec:method}

\begin{figure}[t]
  \centering
  \figplaceholder{\linewidth}{3.2cm}{%
    \textbf{Fig.~3 (pipeline).} Inputs $\gpose$, $\wrench$, $\cpt$
    $\rightarrow$ tip-pose estimator $\rightarrow$ $\tposehat$ $\rightarrow$
    contact-aware steering $\rightarrow$ gripper motion command.}
  \caption{Overview of the approach.}
  \label{fig:pipeline}
\end{figure}

\subsection{Contact as an ad-hoc joint}
\label{sec:method:reach}
\begin{outline}
  \item Model the rod as a quasi-static Kirchhoff rod~\cite{bretl2014quasistatic};
    the edge contact is a point constraint that transmits force but no moment,
    \ie a passive joint whose location along the rod, $s_{\mathrm{c}}$, changes
    as the rod is fed over the edge.
  \item Compute $\Rfree$ and $\Rcon$ by sampling gripper poses in simulation
    (\eg with discrete elastic rods~\cite{bergou2008der}) and validate on
    hardware; this tests \ref{hyp:reach}.
  \item \todo{Analytic result for the planar case (elastica)? Bound on the
    achievable change in tip direction as a function of clearance.}
\end{outline}

\subsection{Sensorized edge}
\label{sec:method:edge}
\begin{outline}
  \item Conductive TPU ring around the opening that reports contact and the
    angular contact location $\cang$ along the ring.
    \todo{Sensing principle, resolution, latency, calibration.}
  \item Design choice: measure the contact location directly instead of
    inferring it from kinematics or wrench (\cf~\cite{bajo2012kinematics}), to
    isolate the value of contact information (\ref{hyp:estimate}).
\end{outline}

\subsection{Haptic tip-pose estimation}
\label{sec:method:estimation}
\begin{outline}
  \item Learn $f_{\theta}:\vect{o}_t\mapsto\tposehat$ from data collected on
    the real system (\cref{sec:experiments}); camera ground truth is used only
    as training labels.
  \item Inputs: $\gpose$, the tared and gravity-compensated wrench $\wrench$,
    the contact state, and $\cang$. Outputs: tip position and direction
    ($(x,y,\theta)$ in the planar case).
  \item Candidate models to compare: MLP; a temporal model over a short history
    (\eg 1D-CNN or GRU) to capture friction and hysteresis; a physics-informed
    model that learns a residual on top of a rod-model
    prediction~\cite{mishani2022realtime}. \todo{Final list of models.}
  \item Loss: weighted position and orientation error; uncertainty from an
    ensemble to gate the controller. \todo{Decide.}
\end{outline}

\subsection{Contact-aware steering}
\label{sec:method:steering}
\begin{outline}
  \item Phases: (i) approach and make contact at a chosen edge location;
    (ii) pivot about the contact to rotate the tip direction; (iii) feed the
    rod to translate the tip to $\tposegoal$.
  \item Closed loop on $\tposehat$: local Jacobian-based servoing (Jacobian
    estimated online or from the learned model), with limits on contact force
    and rod curvature. \todo{Alternative: sampling-based planning in simulation
    with feedback correction.}
\end{outline}

\subsection{Implementation}
\label{sec:method:implementation}
\begin{outline}
  \item ROS~1 (Noetic) with MoveIt for motion planning; synchronized logging of
    joint states, wrench, contact signal, and camera frames.
    \todo{Synchronization method and rates.}
\end{outline}

%% ============================================================================
%% VI. Experimental Plan and Preliminary Results              target: ~1.3 pages
%% Stage 1: experiments may be missing or preliminary, but the plan must show
%% how each hypothesis will be supported or falsified (reviewers score the
%% "feasibility of conclusion"). Stage 2: rename to "Experiments" and report
%% the experiments required by the Extension Charter.
%% Anonymity: no lab names or logos in photos or text.
%% ============================================================================
\section{Experimental Plan and Preliminary Results}
\label{sec:experiments}

\begin{figure}[t]
  \centering
  \figplaceholder{\linewidth}{4cm}{%
    \textbf{Fig.~4 (setup).} Robot arm with wrist force/torque sensor and
    gripper holding the Nitinol rod; sensorized ring (edge); RGB-D camera used
    for ground truth only.}
  \caption{Experimental setup.}
  \label{fig:setup}
\end{figure}

\subsection{Setup}
\label{sec:exp:setup}
\begin{outline}
  \item Yaskawa Motoman SDA10F (one arm), Bota SenseONE EtherCAT force/torque
    sensor at the wrist, parallel gripper \todo{model}, Nitinol rod
    \todo{length, diameter}, conductive-TPU sensorized ring, Intel RealSense
    D435 RGB-D camera for ground truth only \todo{tip marker / tracking method}.
  \item Planar (2D) experiments first, then spatial (3D).
\end{outline}

\subsection{Data collection and protocol}
\label{sec:exp:protocol}
\begin{outline}
  \item Random and structured gripper trajectories that cover contact and
    no-contact states over a range of deformations; each sample is
    $(\gpose,\wrench,\cpt,\tpose)$.
  \item Train/test split by trajectory (testing on unseen configurations);
    repeated trials to quantify sensor noise and hysteresis.
  \item Mitigations for sensor noise, inaccurate contact information, and
    synchronization errors: force/torque tare and calibration, time
    synchronization, camera--robot extrinsic calibration, repeated controlled
    trials.
  \item \todo{Dataset size, number of rods and edges, range of contact
    locations.}
\end{outline}

\subsection{Metrics and baselines}
\label{sec:exp:metrics}
\begin{outline}
  \item Estimation: tip position error [mm] and direction error [deg] on
    held-out data.
  \item Manipulation: success rate within $(\epsilon_p,\epsilon_\theta)$, final
    error, time to target, peak contact force.
  \item Baselines and ablations: inputs ($\gpose$; $\gpose{+}\wrench$;
    $\gpose{+}\wrench{+}\cpt$); model comparison; a rod-model (physics)
    estimator; contact-free and open-loop steering.
\end{outline}

\begin{table}[t]
  \centering
  \caption{How each hypothesis is tested}
  \label{tab:plan}
  \footnotesize
  \begin{tabular}{@{}p{0.06\linewidth}p{0.40\linewidth}p{0.44\linewidth}@{}}
    \toprule
     & Experiment & Supported if / refuted if \\
    \midrule
    \ref{hyp:reach}    & Reachable tip directions with vs.\ without contact
                         (simulation and hardware)
                       & Larger range with contact / no gain \\
    \ref{hyp:estimate} & Estimation error on unseen data; input ablation
                       & Error within tolerance and significant gain from
                         contact input / not \\
    \ref{hyp:steer}    & Steering to targets behind the edge vs.\ baselines
                       & Higher success, lower final error / not \\
    \bottomrule
  \end{tabular}
\end{table}

\subsection{Preliminary results}
\label{sec:exp:prelim}

\begin{table}[t]
  \centering
  \caption{Preliminary tip-pose estimation error (planar, held-out data)}
  \label{tab:prelim}
  \footnotesize
  \begin{tabular}{@{}lcc@{}}
    \toprule
    Inputs & Position [mm] & Direction [deg] \\
    \midrule
    $\gpose$                    & -- & -- \\
    $\gpose+\wrench$            & -- & -- \\
    $\gpose+\wrench+\cpt$       & -- & -- \\
    \bottomrule
  \end{tabular}
\end{table}

\begin{outline}
  \item \todo{First planar dataset: size, error statistics (\cref{tab:prelim}),
    and what they suggest about \ref{hyp:estimate}.}
  \item \todo{Reachability in simulation for \ref{hyp:reach}.}
\end{outline}

\subsection{Planned experiments (Stage 2)}
\label{sec:exp:planned}
\begin{outline}
  \item Spatial (3D) experiments; closed-loop steering trials
    (\ref{hyp:steer}); generalization to other rods and edge geometries.
\end{outline}

%% ============================================================================
%% VII. Conclusion and Expected Insights                     target: ~0.25 page
%% Stage 1 asks for the expected scientific insights, not only results.
%% ============================================================================
\section{Conclusion and Expected Insights}
\label{sec:conclusion}

\begin{outline}
  \item Summary: a contact point with an environment edge, used as an ad-hoc
    joint, together with vision-free tip-pose estimation, for steering the tip
    of an elastic rod in constrained, occluded spaces.
  \item Expected insights: the conditions under which contact enlarges
    reachability; how much tip-pose information the wrench and the contact
    location carry; failure modes (sliding, loss of contact, multiple
    contacts).
  \item Limitations: a single contact, quasi-static motion, known edge
    geometry, one rod material; extension to cables with low bending stiffness
    is open.
  \item Stage-2 plan: 3D experiments and closed-loop steering trials, to be
    refined by the Extension Charter.
\end{outline}


%% Acknowledgments are not allowed in the anonymous submission.
\ifanonymous\else
\section*{Acknowledgments}
\todo{Funding and acknowledgments (camera-ready only).}
\fi

%% Use plainnat to work nicely with natbib (RSS template).
\bibliographystyle{plainnat}
\bibliography{references}

\end{document}
"
%% ---------------------------------------------------------------------------
\providecommand*{\papermode}{draft}
\newcommand*{\paperid}{[add your ID here]} % OpenReview paper ID (Stage 1)
\newcommand*{\papertitle}{Elastic Rod Manipulation Using a Contact Point}

%% ---------------------------------------------------------------------------
%% Packages from the RSS template (hyperref is loaded further down, after the
%% packages it must follow)
%% ---------------------------------------------------------------------------
\usepackage{times}
% numbers option provides compact numerical references in the text.
\usepackage[numbers]{natbib}
\usepackage{multicol}

%% ---------------------------------------------------------------------------
%% Additional packages (none of them changes fonts, margins, or spacing)
%% Avoid: subcaption/caption, enumitem, geometry, mathptmx/newtx, titlesec --
%% they alter the IEEEtran layout.
%% ---------------------------------------------------------------------------
\usepackage{amsmath}
\usepackage{amssymb}
\usepackage{bm}
\usepackage{graphicx}
\graphicspath{{figures/}}
\usepackage[caption=false,font=footnotesize]{subfig} % IEEEtran-safe subfigures
\usepackage{booktabs}
\usepackage{multirow}
\usepackage{siunitx}  % \SI{2}{\milli\metre}, \SI{5}{\degree}
\usepackage{xcolor}
\usepackage{algorithm}
\usepackage{algpseudocode}
\usepackage{xspace}
\usepackage{etoolbox}

\usepackage[bookmarks=true]{hyperref} % RSS template package
\usepackage[capitalise]{cleveref}     % must come after hyperref
\crefname{section}{Sec.}{Secs.}
\crefname{subsection}{Sec.}{Secs.}
\crefname{figure}{Fig.}{Figs.}
\crefname{table}{Table}{Tables}
\crefname{equation}{Eq.}{Eqs.}
\crefname{algorithm}{Alg.}{Algs.}

%% ---------------------------------------------------------------------------
%% Build-mode switches
%% ---------------------------------------------------------------------------
\newif\ifanonymous
\newif\ifshownotes
\ifdefstring{\papermode}{submission}{\anonymoustrue\shownotesfalse}{%
  \ifdefstring{\papermode}{camera}{\anonymousfalse\shownotesfalse}{%
    \anonymousfalse\shownotestrue}}% draft (also the fallback)
\typeout{*** Paper mode: \papermode}

%% PDF metadata (must be anonymous in the review stages)
\hypersetup{
  pdftitle    = {\papertitle},
  pdfkeywords = {deformable linear objects; elastic rods; contact; haptic sensing; manipulation},
}
\ifanonymous
  \hypersetup{pdfauthor={}}
\else
  \hypersetup{pdfauthor={Stav Ashur, Yuval Cassuto, Avishai Sintov}}
\fi

%% ---------------------------------------------------------------------------
%% Draft notes: visible only in draft mode, removed in submission/camera modes
%% (do not also load the todonotes package: it defines \todo as well)
%% ---------------------------------------------------------------------------
\definecolor{notecolor}{RGB}{200,0,0}
\definecolor{outlinecolor}{RGB}{0,70,150}
\DeclareRobustCommand{\draftnote}[3]{%
  \ifshownotes\textcolor{#1}{[\textbf{#2:} #3]}\else\ifhmode\unskip\fi\fi}
\newcommand{\todo}[1]{\draftnote{notecolor}{TODO}{#1}}
\newcommand{\stav}[1]{\draftnote{orange}{Stav}{#1}}
\newcommand{\yuval}[1]{\draftnote{teal}{Yuval}{#1}}
\newcommand{\avishai}[1]{\draftnote{violet}{Avishai}{#1}}

%% Skeleton outlines: bullet lists of planned content. Replace each outline
%% block with prose as you write; they stay blue in every mode so leftovers
%% are easy to spot, and non-draft builds log how many are left.
\newcounter{outlineblocks}
\newenvironment{outline}{%
  \stepcounter{outlineblocks}\par\color{outlinecolor}\begin{itemize}}%
  {\end{itemize}}
\AtEndDocument{\ifshownotes\else\ifnum\value{outlineblocks}>0
  \typeout{WARNING: \arabic{outlineblocks} skeleton outline block(s) left in the paper}%
\fi\fi}

%% \figplaceholder{<width>}{<height>}{<description>}: framed stand-in for a
%% figure that does not exist yet. Replace with \includegraphics{...}.
\newcommand{\figplaceholder}[3]{%
  \fbox{\parbox[c][#2][c]{\dimexpr #1-2\fboxsep-2\fboxrule\relax}{%
    \centering\footnotesize\color{outlinecolor}#3}}}

%% ---------------------------------------------------------------------------
%% Hypotheses and problem statements
%%   \begin{hypothesis}[Short name]\label{hyp:x} ... \end{hypothesis}
%%   \ref{hyp:x} prints "H1", "H2", ...
%% ---------------------------------------------------------------------------
\newcounter{hypothesis}
\renewcommand{\thehypothesis}{H\arabic{hypothesis}}
\crefname{hypothesis}{hypothesis}{hypotheses}
\newenvironment{hypothesis}[1][]{%
  \par\smallskip\refstepcounter{hypothesis}\noindent
  \textbf{\thehypothesis\ifstrempty{#1}{}{ (#1)}:}\hspace{0.5em}\itshape\ignorespaces}%
  {\par\smallskip}
\newtheorem{problem}{Problem}
\crefname{problem}{Problem}{Problems}

%% ---------------------------------------------------------------------------
%% Notation
%% ---------------------------------------------------------------------------
\newcommand{\R}{\mathbb{R}}
\newcommand{\SE}[1]{\mathrm{SE}(#1)}
\newcommand{\SO}[1]{\mathrm{SO}(#1)}
\newcommand{\vect}[1]{\bm{#1}}
\newcommand{\gpose}{\vect{T}_{\mathrm{g}}}          % gripper (end-effector) pose
\newcommand{\tpose}{\vect{T}_{\mathrm{tip}}}        % DLO tip pose
\newcommand{\tposehat}{\hat{\vect{T}}_{\mathrm{tip}}} % estimated tip pose
\newcommand{\tposegoal}{\vect{T}^{*}_{\mathrm{tip}}} % target tip pose
\newcommand{\wrench}{\vect{w}}                      % wrist F/T reading (f, tau)
\newcommand{\cpt}{\vect{c}}                         % contact measurement
\newcommand{\cang}{\phi_{\mathrm{c}}}               % contact angle on the ring
\newcommand{\Rfree}{\mathcal{R}_{\mathrm{free}}}    % reachable set, no contact
\newcommand{\Rcon}{\mathcal{R}_{\mathrm{c}}}        % reachable set, with contact

\newcommand{\eg}{e.g.,\xspace}
\newcommand{\ie}{i.e.,\xspace}
\newcommand{\etal}{et~al.\xspace}
\newcommand{\wrt}{w.r.t.\xspace}
\newcommand{\cf}{cf.\xspace}

\begin{document}

\title{\papertitle}

\ifanonymous
  % Author block required by the RSS template for anonymous review
  \author{Author Names Omitted for Anonymous Review. Paper-ID \paperid}
\else
  \author{\authorblockN{Stav Ashur, Yuval Cassuto, and Avishai Sintov}
  \authorblockA{Tel Aviv University, Tel Aviv, Israel}}
  % TODO (camera-ready): departments and e-mail addresses
\fi

\maketitle

\begin{abstract}
Sensors and devices such as cameras and brushes can be mounted on the tips of
elastic rods and other deformable linear objects (DLOs), allowing them to be
inserted into and manipulated within constrained spaces. Using the right
contact point enables the DLO to reach poses that are infeasible when no
contact with the environment is acceptable. Such a contact point serves as an
ad-hoc joint, facilitating direction changes of the DLO's tip that are
impossible if no collisions are tolerated. In this work, we propose the use of
haptic sensors to perform manipulation that brings the tip of the DLO to a
target pose in a constrained environment by utilizing a contact point of the
DLO with an environment edge.
\todo{Add 1--2 sentences on the hypotheses and methodology (vision-free tip-pose
estimation from gripper pose, wrist force/torque, and contact location;
contact-aware steering) and on the preliminary results.}
\end{abstract}

\IEEEpeerreviewmaketitle

%% Sections, in order
%% ============================================================================
%% I. Introduction                               target: ~0.75 page incl. Fig. 1
%% Motivate the problem, state the gap, preview the hypotheses (Sec. IV) and
%% the methodology (Sec. V), and list the contributions.
%% ============================================================================
\section{Introduction}
\label{sec:intro}

\begin{outline}
  \item \textbf{Motivation.} Cameras, brushes, and other tools are mounted on
    the tips of elastic rods and other DLOs and inserted into constrained
    spaces, \eg for inspection and repair inside engines~\cite{dong2017slender}.
    \todo{1--2 more application examples.}
  \item \textbf{Challenge.} The rod's shape depends on mechanics that are hard
    to model~\cite{sanchez2018survey}, and once inserted, the rod is hidden by
    the environment or by the robot itself, so visual state estimation is
    unreliable.
  \item \textbf{Key idea.} Contact with the environment is usually avoided.
    Instead, we use a contact point between the rod and an environment edge as
    an \emph{ad-hoc joint}: pivoting about it changes the tip direction in ways
    that are impossible if no collisions are tolerated (\cref{fig:teaser}).
  \item \textbf{Problem.} Bring the tip of a gripped elastic rod to a target
    pose in a constrained environment by exploiting a contact with an
    environment edge, using only haptic sensing (robot proprioception, wrist
    force/torque, contact location) during execution (\cref{sec:problem}).
  \item \textbf{Approach.} A sensorized edge reports where the rod touches it;
    a learned estimator maps gripper pose, wrench, and contact location to the
    tip pose; a contact-aware controller steers the tip to the target
    (\cref{sec:method}).
\end{outline}

% Fig. 1 is defined here, not at the top of the section, so that it lands at
% the top of the right column of page 1. LaTeX puts a [t] float at the top of
% the left column (above the abstract) if that column still has room when the
% float is reached; if this happens after edits, move this block further down.
\begin{figure}[t]
  \centering
  \figplaceholder{\linewidth}{4.2cm}{%
    \textbf{Fig.~1 (teaser).} A robot grasps an elastic rod whose tip carries a
    tool. The rod rests on an environment edge (sensorized ring) that acts as an
    ad-hoc joint, letting the tip turn toward a target pose that is unreachable
    without contact. Inset: reachable tip directions with vs.\ without contact.}
  \caption{Using a contact point with an environment edge as an ad-hoc joint
    lets the tip of a gripped elastic rod reach poses that are infeasible
    without contact.}
  \label{fig:teaser}
\end{figure}

\begin{outline}
  \item \textbf{Hypotheses} (\cref{sec:hypotheses}): contact enlarges the set
    of reachable tip poses (\ref{hyp:reach}); haptic signals suffice to estimate
    the tip pose without vision (\ref{hyp:estimate}); closed-loop haptic
    steering reaches targets that contact-free or open-loop strategies miss
    (\ref{hyp:steer}).
  \item \textbf{Contributions}:
    \begin{itemize}
      \item a formulation of DLO tip manipulation that uses a single
        environment contact as an ad-hoc joint, with a reachability analysis;
      \item a vision-free tip-pose estimator that uses proprioception,
        force/torque, and contact location;
      \item an experimental methodology and dataset with a Nitinol rod, from
        planar (2D) to spatial (3D) settings;
      \item preliminary results on the hypotheses.
        \todo{Update after the experiments.}
    \end{itemize}
\end{outline}

%% ============================================================================
%% II. Related Work                                           target: ~0.5 page
%% Stage 1: cite only the work most directly tied to the central problem (what
%% is needed to establish novelty). Stage 2 requires a complete section.
%% Anonymity: cite the authors' own prior work in the third person.
%% RSS 2027 checks references automatically and desk-rejects hallucinated
%% citations, so check every entry you add to references.bib.
%% Merged with the notes in material/2_Related_work.tex (earlier document).
%% ============================================================================
\section{Related Work}
\label{sec:related}

\subsection{DLO state estimation without vision}
\begin{outline}
  \item Vision is the dominant modality for DLO perception and
    manipulation~\cite{sanchez2018survey}, but it fails under the occlusions we
    target. Force/torque (F/T) and tactile sensing are therefore common in DLO
    insertion and
    routing~\cite{kr-aecsagphi-22,nkot-sditfw-97,slk-ftmdloeffi-22,wjly-cratmp-22,zdpp-dhattfmn-19,lvw-crcrisb-24}.
  \item For an elastic rod in static equilibrium, the shape is determined by
    the wrench at one end~\cite{bretl2014quasistatic}; this enabled real-time,
    vision-free shape estimation of a wire held by two arms from F/T
    sensing~\cite{mishani2022realtime}.
  \item \citet{slk-ftmdloeffi-22} manipulate a DLO by keeping it under tension,
    assuming it is straight when tensioned and that gravity makes free segments
    vertical; F/T readings detect the tension and feed an analytic model for
    estimating the DLO state.
  \item Gap: these works assume a contact-free rod held at both ends, or a DLO
    kept taut. An environment contact that bends the rod adds an unknown
    reaction force that breaks the wrench-to-shape mapping; we add the contact
    location to recover it.
\end{outline}

\subsection{DLO insertion and threading}
\begin{outline}
  \item \citet{nkot-sditfw-97} insert a flexible wire using a stereo camera and
    an F/T sensor, first straightening its free end; the insertion strategy is
    standard, and deformation caused by contact with the environment is not
    discussed.
  \item Tight-tolerance threading of string and rope through
    openings~\cite{wang2015online}; learned DLO insertion with flexibility
    estimated from visual cues~\cite{lc-ldloilfevc-24}.
  \item Tactile insertion: \citet{zdpp-dhattfmn-19} learn to guide a deformed
    DLO through a hole using a 16-taxel tactile sensor, without assuming
    compliance or deforming the DLO. \citet{wjly-cratmp-22} build a library of
    tactile-driven primitives (cable following, pivoting, weaving, insertion)
    using a GelSight sensor, designed for the NIST Assembly Task Board; each
    primitive is a state machine whose transitions are triggered by tactile
    input, and insertion is standard: push until the force drops and a rise in
    torque shows that rotation is impossible. Tactile feedback also enables
    following a cable in the grip~\cite{she2021cable}.
  \item \citet{lvw-crcrisb-24} use reinforcement learning (RL) with F/T sensing
    to push compliant objects, including DLOs, into tunnels of unknown shape; a
    Dynamic Safety Lock computes the forces acting on the end effector to allow
    some pushing but stops when the resistance is too high. This is the only
    work we know of with non-straight passages, but its bends do not exceed
    roughly $\pi/6$, and it addresses the uncertainty of the passage shape
    rather than motion planning or control.
  \item Gap: insertion either relies on vision or pushes along a straight or
    mildly bent passage until the force changes; contact with the passage is
    treated as resistance or uncertainty, not used to redirect the tip, so
    angular threading behind an occluding edge is not addressed.
\end{outline}

\subsection{Exploiting environment contacts}
\begin{outline}
  \item Extrinsic dexterity uses external contacts to extend what a simple
    gripper can do~\cite{dafle2014extrinsic}. For peg-in-hole insertion,
    \citet{kr-aecsagphi-22} estimate the contact line between the peg and the
    hole with a GelSight sensor and feed it to an RL model that selects the
    next action.
  \item For DLOs, environmental contacts and fixtures have been used to shape
    cables~\cite{zhu2020environmental} and to manipulate them blindly with
    force feedback~\cite{slk-ftmdloeffi-22}.
  \item Contacts along slender elastic bodies can be detected and localized
    from kinematics, as done for continuum robots~\cite{bajo2012kinematics}.
  \item Gap: contacts are used to shape or hold the DLO body, or are sensed on a
    rigid peg, not used as a joint for steering a \emph{free tip} to a target
    pose with haptic feedback.
    \todo{Confirm this novelty claim against recent work (2024--2026).}
\end{outline}

%% ============================================================================
%% III. Problem Statement                                     target: ~0.5 page
%% Stage 1: the problem statement must be fully developed (a concrete technical
%% formulation, not a high-level argument).
%% ============================================================================
\section{Problem Statement}
\label{sec:problem}

\begin{figure}[t]
  \centering
  \figplaceholder{\linewidth}{3.6cm}{%
    \textbf{Fig.~2 (notation).} Gripper pose $\gpose$, rod of length $L$,
    contact at arc length $s_{\mathrm{c}}$ and ring angle $\cang$, wrist wrench
    $\wrench$, tip pose $\tpose$, and target $\tposegoal$.}
  \caption{Problem setup and notation.}
  \label{fig:notation}
\end{figure}

\begin{outline}
  \item \textbf{Object.} A naturally straight, superelastic Nitinol rod of
    length $L$, modeled as an inextensible Kirchhoff rod under quasi-static
    manipulation~\cite{bretl2014quasistatic}, grasped rigidly at one end; the
    free tip carries a tool. \todo{Rod length, diameter, material constants.}
  \item \textbf{Robot and sensing.} Gripper pose $\gpose\in\SE{3}$ ($\SE{2}$
    in planar experiments) from forward kinematics; wrist wrench
    $\wrench=(\vect{f},\vect{\tau})\in\R^{6}$ from a force/torque sensor; a
    sensorized environment edge (ring) that reports contact and the contact
    angle $\cang$ along the ring (\cref{fig:notation}).
  \item \textbf{Contact.} The rod touches the edge at a single point at an
    unknown arc length $s_{\mathrm{c}}\in(0,L)$; the edge applies an unknown
    reaction force (normal and friction). \todo{Sticking vs.\ sliding contact;
    friction model.}
  \item \textbf{Observation} at time $t$: $\vect{o}_t=(\gpose,\wrench,\cpt)$,
    where $\cpt$ holds the contact state and $\cang$. No vision is used during
    execution; an RGB-D camera provides ground truth only.
  \item \textbf{Reachability.} Let $\Rfree$ and $\Rcon$ be the sets of tip
    poses reachable in the workspace $\mathcal{W}$ without and with an edge
    contact, respectively.
  \item \textbf{Assumptions.} Quasi-static motion, a single contact, known edge
    geometry, deformation within the rod's elastic range.
    \todo{Gravity: modeled or negligible?}
\end{outline}

\begin{problem}[Haptic tip-pose estimation]\label{prob:estimation}
Given the observation $\vect{o}_t$, estimate the tip pose, $\tposehat\approx\tpose$.
\end{problem}

\begin{problem}[Contact-aided tip manipulation]\label{prob:manipulation}
Given a target $\tposegoal\in\Rcon\setminus\Rfree$, find gripper motions that
bring $\tpose$ to within position and orientation tolerances
$(\epsilon_p,\epsilon_\theta)$ of $\tposegoal$ by using a contact with the
edge, while avoiding other collisions and keeping the rod within its elastic
range.
\end{problem}

%% ============================================================================
%% IV. Hypotheses                                             target: ~0.4 page
%% Stage-1 reviewers score (i) the value of the hypotheses and (ii) whether the
%% planned evaluation can support OR falsify them. Make each hypothesis
%% testable and state what result would refute it. Empirical hypotheses are OK.
%% Reference them as \ref{hyp:...} (prints H1, H2, ...).
%% ============================================================================
\section{Hypotheses}
\label{sec:hypotheses}

\todo{Refine the wording and add quantitative thresholds once the setup is fixed.}

\begin{hypothesis}[Reachability]\label{hyp:reach}
Allowing a single contact between the rod and an environment edge strictly
enlarges the set of reachable tip poses, in particular the range of tip
directions, compared with contact-free manipulation in the same constrained
workspace ($\Rfree\subsetneq\Rcon$).
\end{hypothesis}
\begin{outline}
  \item Rationale: the contact acts as an ad-hoc joint (pivot), and its
    location along the rod is controlled by feeding the rod over the edge.
  \item Refuted if, for the tested clearances, the tip directions reachable
    with contact do not exceed those reachable without it (simulation and
    hardware).
\end{outline}

\begin{hypothesis}[Vision-free estimation]\label{hyp:estimate}
Gripper pose, wrist wrench, and contact location suffice to estimate the tip
pose without vision, and the contact location significantly reduces the error
compared with gripper pose and wrench alone.
\end{hypothesis}
\begin{outline}
  \item Rationale: without contact, the wrench at the gripped end determines
    the rod's shape~\cite{bretl2014quasistatic,mishani2022realtime}; a contact
    adds an unknown reaction force, and knowing where it acts recovers the
    missing information.
  \item Refuted if the tip error on unseen configurations exceeds the task
    tolerance, or if adding the contact location gives no significant
    improvement (input ablation, \cref{sec:experiments}).
\end{outline}

\begin{hypothesis}[Contact-aided steering]\label{hyp:steer}
Closed-loop steering with the haptic tip-pose estimate brings the tip to
target poses behind the edge with a higher success rate than contact-free and
open-loop strategies.
\end{hypothesis}
\begin{outline}
  \item Refuted if the success rate and final error are not better than those
    of the baselines on the same targets.
\end{outline}

\begin{outline}
  \item \textbf{Expected insights:} when contact helps and by how much (as a
    function of rod stiffness, clearance, and contact location), and which
    haptic signals carry the tip-pose information.
\end{outline}

%% ============================================================================
%% V. Methodology                                             target: ~1.3 pages
%% Stage 1: the methodology must be fully developed, including the specific
%% design choices in the system architecture and the experiments.
%% ============================================================================
\section{Methodology}
\label{sec:method}

\begin{figure}[t]
  \centering
  \figplaceholder{\linewidth}{3.2cm}{%
    \textbf{Fig.~3 (pipeline).} Inputs $\gpose$, $\wrench$, $\cpt$
    $\rightarrow$ tip-pose estimator $\rightarrow$ $\tposehat$ $\rightarrow$
    contact-aware steering $\rightarrow$ gripper motion command.}
  \caption{Overview of the approach.}
  \label{fig:pipeline}
\end{figure}

\subsection{Contact as an ad-hoc joint}
\label{sec:method:reach}
\begin{outline}
  \item Model the rod as a quasi-static Kirchhoff rod~\cite{bretl2014quasistatic};
    the edge contact is a point constraint that transmits force but no moment,
    \ie a passive joint whose location along the rod, $s_{\mathrm{c}}$, changes
    as the rod is fed over the edge.
  \item Compute $\Rfree$ and $\Rcon$ by sampling gripper poses in simulation
    (\eg with discrete elastic rods~\cite{bergou2008der}) and validate on
    hardware; this tests \ref{hyp:reach}.
  \item \todo{Analytic result for the planar case (elastica)? Bound on the
    achievable change in tip direction as a function of clearance.}
\end{outline}

\subsection{Sensorized edge}
\label{sec:method:edge}
\begin{outline}
  \item Conductive TPU ring around the opening that reports contact and the
    angular contact location $\cang$ along the ring.
    \todo{Sensing principle, resolution, latency, calibration.}
  \item Design choice: measure the contact location directly instead of
    inferring it from kinematics or wrench (\cf~\cite{bajo2012kinematics}), to
    isolate the value of contact information (\ref{hyp:estimate}).
\end{outline}

\subsection{Haptic tip-pose estimation}
\label{sec:method:estimation}
\begin{outline}
  \item Learn $f_{\theta}:\vect{o}_t\mapsto\tposehat$ from data collected on
    the real system (\cref{sec:experiments}); camera ground truth is used only
    as training labels.
  \item Inputs: $\gpose$, the tared and gravity-compensated wrench $\wrench$,
    the contact state, and $\cang$. Outputs: tip position and direction
    ($(x,y,\theta)$ in the planar case).
  \item Candidate models to compare: MLP; a temporal model over a short history
    (\eg 1D-CNN or GRU) to capture friction and hysteresis; a physics-informed
    model that learns a residual on top of a rod-model
    prediction~\cite{mishani2022realtime}. \todo{Final list of models.}
  \item Loss: weighted position and orientation error; uncertainty from an
    ensemble to gate the controller. \todo{Decide.}
\end{outline}

\subsection{Contact-aware steering}
\label{sec:method:steering}
\begin{outline}
  \item Phases: (i) approach and make contact at a chosen edge location;
    (ii) pivot about the contact to rotate the tip direction; (iii) feed the
    rod to translate the tip to $\tposegoal$.
  \item Closed loop on $\tposehat$: local Jacobian-based servoing (Jacobian
    estimated online or from the learned model), with limits on contact force
    and rod curvature. \todo{Alternative: sampling-based planning in simulation
    with feedback correction.}
\end{outline}

\subsection{Implementation}
\label{sec:method:implementation}
\begin{outline}
  \item ROS~1 (Noetic) with MoveIt for motion planning; synchronized logging of
    joint states, wrench, contact signal, and camera frames.
    \todo{Synchronization method and rates.}
\end{outline}

%% ============================================================================
%% VI. Experimental Plan and Preliminary Results              target: ~1.3 pages
%% Stage 1: experiments may be missing or preliminary, but the plan must show
%% how each hypothesis will be supported or falsified (reviewers score the
%% "feasibility of conclusion"). Stage 2: rename to "Experiments" and report
%% the experiments required by the Extension Charter.
%% Anonymity: no lab names or logos in photos or text.
%% ============================================================================
\section{Experimental Plan and Preliminary Results}
\label{sec:experiments}

\begin{figure}[t]
  \centering
  \figplaceholder{\linewidth}{4cm}{%
    \textbf{Fig.~4 (setup).} Robot arm with wrist force/torque sensor and
    gripper holding the Nitinol rod; sensorized ring (edge); RGB-D camera used
    for ground truth only.}
  \caption{Experimental setup.}
  \label{fig:setup}
\end{figure}

\subsection{Setup}
\label{sec:exp:setup}
\begin{outline}
  \item Yaskawa Motoman SDA10F (one arm), Bota SenseONE EtherCAT force/torque
    sensor at the wrist, parallel gripper \todo{model}, Nitinol rod
    \todo{length, diameter}, conductive-TPU sensorized ring, Intel RealSense
    D435 RGB-D camera for ground truth only \todo{tip marker / tracking method}.
  \item Planar (2D) experiments first, then spatial (3D).
\end{outline}

\subsection{Data collection and protocol}
\label{sec:exp:protocol}
\begin{outline}
  \item Random and structured gripper trajectories that cover contact and
    no-contact states over a range of deformations; each sample is
    $(\gpose,\wrench,\cpt,\tpose)$.
  \item Train/test split by trajectory (testing on unseen configurations);
    repeated trials to quantify sensor noise and hysteresis.
  \item Mitigations for sensor noise, inaccurate contact information, and
    synchronization errors: force/torque tare and calibration, time
    synchronization, camera--robot extrinsic calibration, repeated controlled
    trials.
  \item \todo{Dataset size, number of rods and edges, range of contact
    locations.}
\end{outline}

\subsection{Metrics and baselines}
\label{sec:exp:metrics}
\begin{outline}
  \item Estimation: tip position error [mm] and direction error [deg] on
    held-out data.
  \item Manipulation: success rate within $(\epsilon_p,\epsilon_\theta)$, final
    error, time to target, peak contact force.
  \item Baselines and ablations: inputs ($\gpose$; $\gpose{+}\wrench$;
    $\gpose{+}\wrench{+}\cpt$); model comparison; a rod-model (physics)
    estimator; contact-free and open-loop steering.
\end{outline}

\begin{table}[t]
  \centering
  \caption{How each hypothesis is tested}
  \label{tab:plan}
  \footnotesize
  \begin{tabular}{@{}p{0.06\linewidth}p{0.40\linewidth}p{0.44\linewidth}@{}}
    \toprule
     & Experiment & Supported if / refuted if \\
    \midrule
    \ref{hyp:reach}    & Reachable tip directions with vs.\ without contact
                         (simulation and hardware)
                       & Larger range with contact / no gain \\
    \ref{hyp:estimate} & Estimation error on unseen data; input ablation
                       & Error within tolerance and significant gain from
                         contact input / not \\
    \ref{hyp:steer}    & Steering to targets behind the edge vs.\ baselines
                       & Higher success, lower final error / not \\
    \bottomrule
  \end{tabular}
\end{table}

\subsection{Preliminary results}
\label{sec:exp:prelim}

\begin{table}[t]
  \centering
  \caption{Preliminary tip-pose estimation error (planar, held-out data)}
  \label{tab:prelim}
  \footnotesize
  \begin{tabular}{@{}lcc@{}}
    \toprule
    Inputs & Position [mm] & Direction [deg] \\
    \midrule
    $\gpose$                    & -- & -- \\
    $\gpose+\wrench$            & -- & -- \\
    $\gpose+\wrench+\cpt$       & -- & -- \\
    \bottomrule
  \end{tabular}
\end{table}

\begin{outline}
  \item \todo{First planar dataset: size, error statistics (\cref{tab:prelim}),
    and what they suggest about \ref{hyp:estimate}.}
  \item \todo{Reachability in simulation for \ref{hyp:reach}.}
\end{outline}

\subsection{Planned experiments (Stage 2)}
\label{sec:exp:planned}
\begin{outline}
  \item Spatial (3D) experiments; closed-loop steering trials
    (\ref{hyp:steer}); generalization to other rods and edge geometries.
\end{outline}

%% ============================================================================
%% VII. Conclusion and Expected Insights                     target: ~0.25 page
%% Stage 1 asks for the expected scientific insights, not only results.
%% ============================================================================
\section{Conclusion and Expected Insights}
\label{sec:conclusion}

\begin{outline}
  \item Summary: a contact point with an environment edge, used as an ad-hoc
    joint, together with vision-free tip-pose estimation, for steering the tip
    of an elastic rod in constrained, occluded spaces.
  \item Expected insights: the conditions under which contact enlarges
    reachability; how much tip-pose information the wrench and the contact
    location carry; failure modes (sliding, loss of contact, multiple
    contacts).
  \item Limitations: a single contact, quasi-static motion, known edge
    geometry, one rod material; extension to cables with low bending stiffness
    is open.
  \item Stage-2 plan: 3D experiments and closed-loop steering trials, to be
    refined by the Extension Charter.
\end{outline}


%% Acknowledgments are not allowed in the anonymous submission.
\ifanonymous\else
\section*{Acknowledgments}
\todo{Funding and acknowledgments (camera-ready only).}
\fi

%% Use plainnat to work nicely with natbib (RSS template).
\bibliographystyle{plainnat}
\bibliography{references}

\end{document}
